%
% CHAPTER: Machine Learning
%

\chapterimage{Capek_play.pdf}

\chapter{mnplib}
\label{chap:mnplib}

\begin{quote}
\begin{flushright}
\emph{There are no difficult problems,\\
only lack of imagination.}\\
Antonio García \\
\end{flushright}
\end{quote}
\bigskip

We have seen that the most challenging problems to which we can apply the results of the theory of nescience arise when the set of entities $\mathcal{E}$ under study is composed of abstract elements. The difficulty with abstract entities is that there is no straightforward way to encode them as finite strings of symbols from which they can be effectively reconstructed. Without such an encoding, it becomes impossible to measure their redundancy, accuracy, or any other property required by the theory of nescience.

In practice, a common approach to overcome this limitation is to design an experiment and collect observable outcomes, as is typically done in physics or other empirical sciences. Another possible approach is to collect a set of indirect measurements that serve as observable proxies for the abstract entities, for example, by recording the behavior of users in an online social network to study social interaction patterns.

This chapter is devoted to applying the concept of minimum nescience to the field of machine learning. We assume that the entities under study have been encoded as a dataset $\mathbb{X}$ consisting of $n$ training vectors, each with $p$ predictors (or features), and an associated response variable $\mathbf{y}$ (see Section \ref{sec:machine_learning}).

We will begin by presenting practical approximations for the concepts of miscoding, inaccuracy, and surfeit when the entities are represented as datasets. We will then show how these quantities can be combined into the single measure of nescience. These approximations will allow us to introduce the minimum nescience principle, a method designed to automate the process of selecting optimal models in machine learning (a form of automated machine learning, or AutoML).

In addition to introducing these approximations, we will demonstrate how to apply them to solve practical machine learning problems. The examples provided will be based on the \texttt{nescience}\footnote{\texttt{https://github.com/rleiva/nescience}} library, an open-source Python library that implements the ideas presented in this chapter.

%
% Section: Nescience Library
%

\section{Nescience Python Library}
\label{sec:nescience_library}

The \texttt{nescience} library is an open-source Python package that implements the ideas presented in this book, specifically their application to the field of machine learning. The library follows the API design and conventions of the widely used \texttt{scikit-learn} toolkit, which allows it to be seamlessly combined with the methods and workflows provided by that package.

The \texttt{nescience} library can be installed with the \texttt{pip} utility:

\begin{sourcecode}
{\scriptsize \begin{verbatim}
pip install nescience
\end{verbatim}}
\end{sourcecode}

On the companion website for the library, the reader will find a collection of \texttt{jupyter-lab} notebooks that illustrate how the library works. Each subsection of this chapter is accompanied by a notebook containing implementations of all the examples discussed, so that the reader can reproduce the results and experiment with them interactively. Additional information about the \texttt{nescience} library, including a complete API reference, is also available on the library's website.

%
% Section: The Utility Interface
%

\section{The Utility Interface}
\label{sec:utility_interface}

The practical computation of the components of nescience requires a finite and computable way of measuring the amount of information contained in observed data. The theory is formulated in terms of description length and, ultimately, Kolmogorov complexity, which cannot be computed in general. The \texttt{mnplib} library therefore uses empirical code lengths as practical approximations. The tools to discretize and compress vectors are provided by the Utility Interface \texttt{utils}.

The basic idea is simple. An attribute is represented by a vector of observed values. Equal values are treated as instances of the same symbol, the frequency of each symbol is counted, and the total code length is computed from those frequencies using the ideal coding scheme introduced in Section~\ref{sec:Optimal-Codes}. Numeric attributes require an additional step: before their frequencies can be counted, their values are discretized according to the procedure described in Section~\ref{sec:discretization_continuous_variables}.

No actual compressed bit stream is produced. The resulting value is the ideal number of bits required to encode the observations according to their empirical frequencies.

\subsection*{Encoding Categorical Attributes}

Let $\mathbf{x}=(x_1,\ldots,x_n)$ be an attribute containing $n$ observations, and suppose that the distinct values appearing in the attribute are $g_1,\ldots,g_q$. Let $n_i$ denote the number of times that value $g_i$ occurs in the attribute, so that $\sum_{i=1}^{q} n_i=n$. The relative frequency of $g_i$ is $f_i=\frac{n_i}{n}$.

As discussed in Section~\ref{sec:Optimal-Codes}, a value with relative frequency $f_i$ is assigned the ideal code length $-\log_2 f_i$. Since the value occurs $n_i$ times, the contribution of that value to the code length of the complete attribute is $-n_i\log_2 f_i$. The empirical code length used by \texttt{mnplib} is therefore
$$
\hat{K}_C(\mathbf{x})
=
-\sum_{i:n_i>0}
n_i\log_2\left(\frac{n_i}{n}\right).
$$
This quantity depends only on the frequencies of the distinct values. The particular labels used to represent them are irrelevant. For example, replacing the categories $\{\texttt{red},\texttt{green},\texttt{blue}\}$ with $\{0,1,2\}$ does not modify their frequencies and therefore does not modify the resulting code length. For this reason, \texttt{mnplib} internally replaces categorical values with integer labels. The labels are assigned according to their order of first appearance, but their numerical values have no meaning. They serve only to identify which observations belong to the same category.

Missing categorical values are rejected because they would otherwise introduce ambiguity about whether the missing value represents an actual category or the absence of an observation.

\subsection*{Joint Encoding of Several Attributes}

The same procedure can be applied to several attributes simultaneously. Consider $d$ attributes $\mathbf{x}^{(1)},\ldots,\mathbf{x}^{(d)}$, each containing $n$ observations. For observation $j$, the values of all attributes are grouped into the tuple $\left(x^{(1)}_j,\ldots,x^{(d)}_j\right)$. Each distinct tuple is treated as a joint state. Only joint states that actually occur in the data are considered. The library does not construct the complete Cartesian product of all possible attribute values.

\begin{example}
Given two categorical attributes,
$$
\begin{array}{c|c}
\text{Color} & \text{Shape}\\
\hline
\text{red} & \text{square}\\
\text{red} & \text{circle}\\
\text{blue} & \text{square}\\
\text{red} & \text{square}
\end{array}
$$
the observed joint states are $(\text{red},\text{square}), (\text{red},\text{circle}), (\text{blue},\text{square})$, with frequencies $2$, $1$, and $1$, respectively.
\end{example}

If the observed joint states are $s_1,\ldots,s_q$, with frequencies $n_1,\ldots,n_q$, their code length is computed exactly as for a single attribute:
$$
\hat{K}_C \left(\mathbf{x}^{(1)},\ldots,\mathbf{x}^{(d)} \right) =
-\sum_{i:n_i>0} n_i \log_2 \left( \frac{n_i}{n} \right).
$$

Joint code lengths are particularly important in \texttt{mnplib}, because several components of nescience are computed by comparing the code lengths of individual attributes, subsets of attributes, targets, predictions, and their joint representations.

\subsection*{Encoding Numeric Attributes}

A numeric attribute may contain many different observed values. If every distinct value were treated as a separate symbol, the resulting representation would often contain almost as many states as observations and would not provide a useful description of the regularities in the data. Numeric attributes are therefore discretized before their frequencies are computed.

Let $\mathbf{x}=(x_1,\ldots,x_n)$ be a numeric attribute. The discretization procedure described in Section~\ref{sec:discretization_continuous_variables} divides the interval between the minimum and maximum observed values into $b$ intervals of equal width. Each observation is replaced by the integer label of the interval containing it:
$$
\mathbf{x} \; \longrightarrow\; \tilde{\mathbf{x}} = (\tilde{x}_1,\ldots,\tilde{x}_n).
$$
The resulting labels form a finite set of states. Their frequencies are then counted and the code length is computed in exactly the same way as for a categorical attribute:
$$
\hat{K}_C(\mathbf{x}) \approx \hat{K}_C(\tilde{\mathbf{x}}) =
-\sum_{i:n_i>0} n_i \log_2 \left(\frac{n_i}{n}\right).
$$
Discretization therefore determines the resolution at which numeric values are considered different. Values falling in the same interval are treated as equivalent for the purpose of computing their code length.

\subsection*{Uniform Discretization}

The current implementation of \texttt{mnplib} uses equal-width, or uniform, discretization. Given $b$ bins, the interval between the smallest and largest observed values is divided into $b$ intervals of equal width. Each numeric attribute is discretized independently, so different attributes generally have different interval boundaries even when the same number of bins is used.

The use of equal-width intervals is deliberate. The objective is to preserve, at the selected resolution, differences in the density of observed values.

\begin{example}
Suppose an attribute in which one part of the range contains many observations while another contains only a few. Under equal-width discretization this difference remains visible in the frequencies of the corresponding bins. Those frequencies subsequently affect the code length.

By contrast, an equal-frequency discretization deliberately places approximately the same number of observations in every bin. If all $b$ bins contained exactly $n/b$ observations, their relative frequencies would all be $1/b$, and the resulting code length would be
$$
-\sum_{i=1}^{b} \frac{n}{b} \log_2\left(\frac{1}{b}\right) = n\log_2 b.
$$
The result would then depend almost entirely on the number of bins and would contain little information about the original distribution of values.
\end{example}

For the computation of nescience, the behaviour of equal-frequency is undesirable. Differences in the frequencies of observed values are precisely what allow the code length to distinguish structured attributes from more uniformly distributed ones. Uniform discretization preserves these differences instead of deliberately removing them.

When all observations of a numeric attribute have the same value, no actual discretization is necessary. \texttt{mnplib} assigns all observations to a single state. Its frequency is one and the resulting code length is zero.

\subsection*{Number of Bins}

The number of bins determines the resolution of the representation. Too few bins merge values that may contain relevant differences. Too many bins produce a large number of sparsely populated states and make the observed frequencies increasingly dependent on the particular sample. Section~\ref{sec:discretization_continuous_variables} discusses this trade-off and introduces the two automatic policies implemented by \texttt{mnplib}.

For a single numeric attribute, the default policy is
$$
b_{\mathrm{auto}} = \max \left\{2, \left \lfloor 2n^{1/3} \right \rfloor \right\}.
$$
This rule is inspired by the Rice rule and increases the resolution as more observations become available. The increase is deliberately slower than the increase in sample size, so that both the number of bins and the average number of observations represented by each bin tend to increase with $n$. The minimum value of two prevents a nonconstant numeric attribute from being represented using fewer than two intervals.

The user may also provide an explicit integer number of bins. In that case the supplied value is used directly and does not depend on the number of observations or attributes.

\subsection*{Adaptive Discretization of Several Attributes}

When several attributes are considered jointly, the number of possible combinations of their states grows rapidly with the number of attributes. Even if each attribute contains only a moderate number of states, their joint representation may become sparse. For this reason, \texttt{mnplib} provides an adaptive discretization policy for multidimensional empirical distributions.

For $n$ observations and $d$ jointly supplied attributes, the number of bins assigned to each numeric attribute is
$$
b_{\mathrm{adaptive}} =
\max \left\{2, \left\lfloor \frac{2n^{1/3}}{\log_2(d+1)} \right\rfloor \right\}.
$$
For a single attribute, $d=1$ and therefore $\log_2(d+1)=1$, so the adaptive rule reduces exactly to the ordinary automatic rule. As the number of jointly represented attributes increases, the number of bins assigned to each numeric attribute decreases. The purpose is to compensate, at least partially, for the increasing number of possible joint states and to reduce excessive sparsity. The value $d$ refers to the number of attributes participating in the joint representation, including categorical attributes. Categorical attributes are not themselves discretized, but they contribute additional dimensions to the joint state space and therefore influence the adaptive resolution.

All numeric attributes participating in the same empirical distribution use the same resolved number of bins. Their interval boundaries, however, are computed independently from their own observed minimum and maximum values.

\subsection*{Marginal and Joint Code Lengths}

Care is required when code lengths of different representations are compared\footnote{The implementation Utility Interface is intentionally stateless. Discretization boundaries, encoded values, frequencies, and code lengths are computed from the values supplied to each call. No fitted discretizer or empirical distribution is retained between calls.}. Suppose, for example, that we want to compare the code length of a numeric attribute $\mathbf{x}$ with the joint code length of $\mathbf{x}$ and another attribute $\mathbf{y}$: $\hat{K}_C(\mathbf{x})$, and $\hat{K}_C(\mathbf{x},\mathbf{y})$. The representation of $\mathbf{x}$ should be the same in both computations. Otherwise, part of the observed difference in code length could be caused by using different discretizations rather than by adding $\mathbf{y}$.

This becomes particularly relevant when adaptive discretization is used. A one-dimensional call and a multidimensional call have different values of $d$ and may therefore select different numbers of bins. Consequently, when marginal and joint code lengths are intended to be compared directly, \texttt{mnplib} resolves an appropriate bin count and uses the same explicit integer value for the shared numeric attributes involved in the comparison. Using the same number of bins for the same observed numeric vector produces the same discretization because its minimum and maximum values are unchanged. Different attributes still have their own interval boundaries, since each attribute is discretized independently.

This consistency requirement is important for quantities based on differences or ratios of code lengths, including those used in the computation of miscoding, inaccuracy, and related components of nescience.

%
% Section: Miscoding
%

\section{Miscoding}
\label{sec:machine_learning_miscoding}

In Definition~\ref{def:relative_miscoding} we introduced miscoding as a measure of how inadequately a representation encodes an entity. Its theoretical definition is based on conditional Kolmogorov complexity and therefore cannot, in general, be computed. The \texttt{mnplib} library implements a practical supervised approximation for datasets in which a collection of features $\mathbf{X}$ is used as a representation of a target attribute $\mathbf{y}$.

This setting corresponds to the labeled miscoding problem introduced earlier in Chapter \ref{chap:Labeled_Miscoding}. Each row of the dataset is associated with a target value, and the objective is to determine how much information about the target is missing from the available features, and how much information contained in those features is not recoverable from the target.

The implementation replaces Kolmogorov complexity with the empirical code lengths described in Section~\ref{sec:utility_interface}. For an attribute $\mathbf{x}$, we denote its empirical code length by $\hat K_C(\mathbf{x})$, and the empirical code length of the jointly observed attributes $\mathbf{x}$ and $\mathbf{y}$ by $\hat K_C(\mathbf{x},\mathbf{y})$.

Categorical attributes are encoded directly from the frequencies of their observed states. Numeric attributes are first discretized using the procedure described in Section~\ref{sec:utility_interface}. Consequently, all the quantities introduced in this section are finite and computable approximations of their algorithmic counterparts.

The practical formulation is based on the approximations
\[
K(\mathbf{y}\mid\mathbf{x}) \approx \hat K_C(\mathbf{x},\mathbf{y}) - \hat K_C(\mathbf{x}),
\]
and
\[
K(\mathbf{x}\mid\mathbf{y}) \approx \hat K_C(\mathbf{x},\mathbf{y}) - \hat K_C(\mathbf{y}).
\]
These two quantities lead naturally to the notions of deficiency and surplus.

% Feature level miscoding

\subsection*{Feature-Level Miscoding}

Let $\mathbf{X} = \left[\mathbf{x}_1,\ldots,\mathbf{x}_p\right]$ be a dataset composed of $p$ features and let $\mathbf{y}$ be its target attribute. For each individual feature $\mathbf{x}_j$, \texttt{mnplib} computes three
quantities: the empirical \emph{deficiency} of $\mathbf{x}_j$ with respect to
$\mathbf{y}$ is
\[
\hat d_j = \frac{\hat K_C(\mathbf{x}_j,\mathbf{y}) - \hat K_C(\mathbf{x}_j)}{\hat K_C(\mathbf{y})}.
\]
The empirical \emph{surplus} is
\[
\hat s_j = \frac{\hat K_C(\mathbf{x}_j,\mathbf{y}) - \hat K_C(\mathbf{y})}{\hat K_C(\mathbf{x}_j)}.
\]
Finally, the empirical \emph{miscoding} of the feature is
\[
\hat\mu_j = \max\{\hat d_j,\hat s_j\}.
\]

Deficiency approximates the fraction of the target information that still has to be supplied when the feature is known. A low deficiency therefore indicates that the feature contains most of the information required to reconstruct the target. Surplus measures the amount of information contained in the feature that is not relevant to describing the target. A feature may therefore have low deficiency but substantial surplus: it may contain all the information required to reconstruct the target while also carrying additional information that does not contribute to that description. Miscoding takes the largest of these two defects. A feature has low miscoding only when both directions are adequately represented.

In the implementation, these values are constrained to the interval $[0,1]$. Feature-level quantities are computed during the call to \texttt{fit()} and can subsequently be obtained through

\begin{sourcecode}
{\scriptsize \begin{verbatim}
msd.deficiency_feature()
msd.surplus_feature()
msd.miscoding_feature()
\end{verbatim}}
\end{sourcecode}

or together in tabular form with

\begin{sourcecode}
{\scriptsize \begin{verbatim}
msd.feature_analysis()
\end{verbatim}}
\end{sourcecode}

The latter returns the feature index, feature name, whether the feature was treated as numeric, its empirical code length, and its deficiency, surplus, and miscoding.

\begin{example}
\label{example:miscoding_synthetic}

Consider a synthetic target composed of two independent categorical components $A$ and $B$. We construct five different features:

\begin{itemize}
    \item \texttt{perfect} contains exactly $(A,B)$;
    \item \texttt{deficient} contains only $A$ and therefore lacks $B$;
    \item \texttt{surplus} contains $(A,B,C)$, where $C$ is independent;
    \item \texttt{irrelevant} is independent of the target.
\end{itemize}

\begin{sourcecode}
{\scriptsize \begin{verbatim}
import numpy as np
import pandas as pd

from sklearn.utils import check_random_state
from mnplib.miscoding import Miscoding

rng = check_random_state(42)

A = rng.randint(0, 4, size=6000)
B = rng.randint(0, 4, size=6000)
C = rng.randint(0, 4, size=6000)
D = rng.randint(0, 16, size=6000)

Y = np.array([f"{a}_{b}" for a, b in zip(A, B)])

perfect = np.array([f"{a}_{b}" for a, b in zip(A, B)])

X = pd.DataFrame({
    "perfect"    : perfect,
    "deficient"  : A.astype(str),
    "surplus"    : np.array(
        [f"{a}_{b}_{c}" for a, b, c in zip(A, B, C)]
    ),
    "irrelevant" : D.astype(str)
})

msd = Miscoding().fit(X, Y)

msd.feature_analysis()
\end{verbatim}}
\end{sourcecode}

A representative result is

\begin{center}
\begin{tabular}{lrrr}
\hline
Feature & Deficiency & Surplus & Miscoding \\
\hline
perfect    & 0.000 & 0.000 & 0.000 \\
surplus    & 0.000 & 0.333 & 0.333 \\
deficient  & 0.500 & 0.000 & 0.500 \\
irrelevant & 0.992 & 0.992 & 0.992 \\
\hline
\end{tabular}
\end{center}

The example illustrates the different meanings of the two components. The deficient feature contains only half of the information required to reconstruct the target, but it introduces essentially no surplus. The \texttt{surplus} feature contains all target information, but approximately one third of its own information is not related to the target. The irrelevant feature is poor in both directions.

\end{example}

Redundancy should not be confused with surplus. Consider a feature that contains many repeated or highly similar samples. These samples are redundant because they do not add much new information, but they are not necessarily irrelevant to the target and therefore should not be penalized as surplus. Such redundancy may even be useful in practice, for example by providing repeated evidence that helps some learning algorithms, such as neural networks, estimate patterns more robustly. Surplus is intended to capture information that is unrelated to describing the target, not information that is merely repeated.

The same interface can be used without knowing in advance which features are useful.

\begin{example}
Consider the Iris dataset included in \texttt{scikit-learn}.

\begin{sourcecode}
{\scriptsize \begin{verbatim}
from sklearn.datasets import load_iris
from mnplib.miscoding import Miscoding

iris = load_iris(as_frame=True)
X = iris.data
y = iris.target

msd = Miscoding().fit(X, y)

msd.feature_analysis()
\end{verbatim}}
\end{sourcecode}

A representative result is

\begin{center}
\begin{tabular}{lrrr}
\hline
Feature & Deficiency & Surplus & Miscoding \\
\hline
petal width (cm)  & 0.109 & 0.506 & 0.506 \\
petal length (cm) & 0.148 & 0.525 & 0.525 \\
sepal length (cm) & 0.543 & 0.767 & 0.767 \\
sepal width (cm)  & 0.725 & 0.847 & 0.847 \\
\hline
\end{tabular}
\end{center}

According to miscoding, features "petal width" and "petal length" provide the best individual representations of the target. The important point is not only their final miscoding value. Their separate deficiency and surplus values show why they are preferable.
\end{example}

Individual feature miscoding does not reveal whether two useful features encode essentially the same information. The library therefore also computes a symmetric pairwise miscoding measure, defined for any two features $\mathbf{x}_i$ and $\mathbf{x}_j$ as
\[
\hat\mu(\mathbf{x}_i,\mathbf{x}_j) = \frac{ \hat K_C(\mathbf{x}_i,\mathbf{x}_j)
- \min\{\hat K_C(\mathbf{x}_i), \hat K_C(\mathbf{x}_j) \} }{
\max\{\hat K_C(\mathbf{x}_i), \hat K_C(\mathbf{x}_j)\}}.
\]
Thus, $\hat\mu(\mathbf{x}_i,\mathbf{x}_j) \approx 1$ indicates indicates little common information, whereas $\hat\mu(\mathbf{x}_i,\mathbf{x}_j) \approx 0$ that the two attributes contain largely the same information. The complete matrix can be obtained with

\begin{sourcecode}
{\scriptsize \begin{verbatim}
msd.pairwise_miscoding_matrix()
\end{verbatim}}
\end{sourcecode}

For the Iris dataset the smallest off-diagonal value is obtained for \texttt{petal length} and \texttt{petal width}, with a pairwise miscoding of approximately $0.84$. This is consistent with the fact that these two petal measurements contain a little amount of overlapping information.

\begin{example}

A useful property of empirical code-length miscoding is that it is not restricted to linear relations. Consider $y=x^2$,  with $x$ uniformly distributed over $[-1,1]$. Pearson correlation is close to zero because positive and negative values cancel in the covariance, despite the fact that $x$ determines $y$ exactly.

\begin{sourcecode}
{\scriptsize \begin{verbatim}
from sklearn.feature_selection import r_regression
from sklearn.utils import check_random_state
from mnplib.miscoding import Miscoding

X = check_random_state(42).uniform(
    -1.0, 1.0, size=(5000, 1)
)
y = X[:, 0] ** 2

correlation = r_regression(X, y)[0]

msd = Miscoding().fit(X, y)
miscoding = msd.miscoding_feature(0)

print("Pearson correlation:", correlation)
print("Miscoding:", miscoding)
\end{verbatim}}
\end{sourcecode}

A representative result is

\begin{sourcecode}
{\scriptsize \begin{verbatim}
Pearson correlation: -0.0201
Miscoding:            0.2993
\end{verbatim}}
\end{sourcecode}

Pearson correlation suggests almost no linear relation. Miscoding, in contrast, detects that the two attributes are strongly related.
\end{example}

The miscoding is not exactly zero. Although $x$ completely determines $y$, the converse transformation is not unique: $x=\pm\sqrt{y}$. Consequently, additional information is needed to reconstruct the actual sequence of signs of $x$ from $y$. This contributes to surplus.

Mutual information provides a more general measure of dependence than correlation and is also capable of detecting nonlinear relations. However, it produces a single quantity measuring the information shared by two variables. Miscoding separates the remaining information into two directional components. The distinction can be illustrated with a simple discrete example.

\begin{example}
Let $\mathbf{y}=(A,B)$, where $A$ and $B$ are independent balanced binary attributes. Let the representation be $\mathbf{x}=(A,Z_1,Z_2,Z_3)$, where $Z_1,Z_2,Z_3$ are independent balanced binary attributes unrelated to the target. The feature shares the component $A$ with the target, but it lacks $B$ and contains three additional independent components.

\begin{sourcecode}
{\scriptsize \begin{verbatim}
import itertools
import numpy as np

from sklearn.metrics import mutual_info_score
from mnplib.miscoding import Miscoding

rows = np.array(
    list(itertools.product([0, 1], repeat=5)),
    dtype=int,
)
rows = np.tile(rows, (32, 1))

a, b, z1, z2, z3 = rows.T

X = 8*a + 4*z1 + 2*z2 + z3
y = 2*a + b

mi = mutual_info_score(X, y) / np.log(2)

msd = Miscoding(
    X_type="categorical",
    y_type="categorical",
).fit(X[:, None], y)

deficiency = msd.deficiency_feature(0)
surplus    = msd.surplus_feature(0)

print("Mutual information:", mi)
print("Deficiency:", deficiency)
print("Surplus:", surplus)
\end{verbatim}}
\end{sourcecode}

The result is

\begin{sourcecode}
{\scriptsize \begin{verbatim}
Mutual information: 1.00
Deficiency:         0.50
Surplus:            0.75
\end{verbatim}}
\end{sourcecode}

The mutual information correctly reports that $\mathbf{x}$ and $\mathbf{y}$ share one bit. Miscoding provides additional diagnostic information. The deficiency $\hat d=0.5$ shows that half of the information required to reconstruct the target is missing from $\mathbf{x}$, namely the bit $B$. At the same time, $\hat s=0.75$ shows that three quarters of the information in $\mathbf{x}$ must still be supplied when $\mathbf{y}$ is known, corresponding here to the independent attributes $Z_1,Z_2,Z_3$.

Thus mutual information tells us \emph{how much} information the two attributes share, whereas the decomposition used by miscoding also tells us \emph{how the representation fails}.
\end{example}

It is important to distinguish this conceptual comparison from the interface provided by \texttt{scikit-learn}. Mutual information itself is mathematically defined between a multidimensional vector and a target. However, \texttt{mutual\_info\_classif()} and \texttt{mutual\_info\_regression()} compute one mutual-information value for each individual feature. They do not directly return $I((X_1,\ldots,X_m);Y)$ for an arbitrary selected subset. For discrete data, a joint state may be constructed explicitly, as in the example above; continuous multivariate estimation requires a dedicated estimator. In contrast, as it is shown in the next section, \texttt{Miscoding.subset\_analysis()} directly evaluates selected feature subsets.

\subsection*{Subset-level Miscoding}

Feature-level diagnostics are insufficient when several features are used
jointly. Two individually useful features may contain almost the same
information, while two individually weak features may become highly
informative when considered together.

Let $S\subseteq\{1,\ldots,p\}$ be a subset of feature indices and let $\mathbf{X}_S$ denote the joint attribute formed by the selected features. The empirical deficiency of a feature subset $S$ with respect to the target
$\mathbf{y}$ is
\[
\hat d(S) = \frac{\hat K_C(\mathbf{X}_S,\mathbf{y}) - \hat K_C(\mathbf{X}_S)}{\hat K_C(\mathbf{y})}.
\]
Its empirical surplus is
\[
\hat s(S) = \frac{\hat K_C(\mathbf{X}_S,\mathbf{y}) - \hat K_C(\mathbf{y})}{\hat K_C(\mathbf{X}_S)},
\]
and its empirical miscoding is
\[
\hat\mu(S) = \max\{\hat d(S),\hat s(S)\}.
\]
These quantities are computed directly from the joint empirical distribution. They are not obtained by adding or averaging the individual feature miscodings. This distinction is essential. The information supplied by several attributes depends on their joint structure. In particular, duplicated or strongly redundant features need not improve the representation even when each of them has a low individual miscoding.

Subset diagnostics are available through

\begin{sourcecode}
{\scriptsize \begin{verbatim}
msd.subset_analysis([0, 2, 3])

msd.deficiency_subset([0, 2, 3])
msd.surplus_subset([0, 2, 3])
msd.miscoding_subset([0, 2, 3])
\end{verbatim}}
\end{sourcecode}

For numeric attributes, all marginal and joint code lengths involved in a particular subset calculation use the same resolved number of bins. This is important because otherwise $\hat K_C(\mathbf{X}_S)$, $\hat K_C(\mathbf{y})$, $\hat K_C(\mathbf{X}_S,\mathbf{y})$ could refer to different discretization resolutions and their differences would not be directly comparable. The number of bins is resolved adaptively as described in Section \ref{sec:utility_interface}.

The number of possible joint states grows rapidly as additional features are introduced. With a finite sample, a joint empirical distribution may eventually become so sparse that its code lengths are no longer considered sufficiently reliable. For this reason, \texttt{mnplib} evaluates the observed joint states before returning subset deficiency, surplus, and miscoding.

Let $n$ be the number of observations and $q$ the number of distinct joint states that actually occur. The mean occupancy is $\bar o = n /q$. Let $q_1$ be the number of observed joint states that occur exactly once. The
singleton fraction is $f_1 = q_1 / q$ . The current implementation regards the empirical joint distribution as
reliable when $\bar o \geq 2$ and $f_1\leq 0.5$. If the joint distribution does not satisfy the reliability condition, the subset deficiency, surplus, and miscoding are returned as \texttt{NaN} rather than presenting an unstable estimate as a meaningful result.

The reliability check is deliberately applied to the joint representation $(\mathbf{X}_S,\mathbf{y})$, because this is the distribution on which the conditional code-length approximations ultimately depend.

\subsection*{Feature Selection}

The subset formulation makes it possible to use miscoding directly for feature selection. The method \texttt{select\_features()} implements greedy forward selection. Starting with the empty subset, it evaluates every possible one-feature extension and selects the candidate that produces the smallest reliable subset miscoding. A feature is accepted only if it reduces the current miscoding by more than the requested value of \texttt{min\_improvement}. Selection stops when no reliable candidate remains, when no candidate provides the required improvement, or when \texttt{max\_features} is reached.

Although pairwise miscoding is available separately through \texttt{pairwise\_miscoding\_matrix()}, feature selection does not need to add a separate miscoding penalty. Miscoding is already reflected in the joint code
length of the selected subset: adding a feature that merely repeats information already present in the subset provides little or no reduction in deficiency and may increase surplus. This distinguishes subset miscoding from feature-selection methods that rank features independently.

Strict feature selection is not always appropriate. When constructing a model, we may want an ordering of the available features even after additional features stop reducing subset miscoding. The method \texttt{rank\_features()} therefore follows a related but different greedy procedure. At every step it evaluates the reliable one-feature extensions of the current subset and adds the best candidate. Unlike \texttt{select\_features()}, it does not require the newly added feature to reduce the current miscoding. Ranking continues until the requested number of features has been reached or no reliable extension remains.

The distinction is useful in automatic model construction. Feature selection answers the question "Which features can be justified by an actual reduction of miscoding?" whereas feature ranking answers "In which order should additional features be considered?" The latter is particularly useful when another criterion, such as inaccuracy, surfeit, or total nescience, will ultimately decide how many of the ranked features should be used by the model.

\begin{example}
Let $A$ and $B$ be two binary attributes, and let $Y=2A+B$. The dataset contains $A$, an exact copy of $A$, and $B$.

\begin{sourcecode}
{\scriptsize \begin{verbatim}
import numpy as np

from functools import partial
from sklearn.feature_selection import (
    SelectKBest,
    mutual_info_classif,
)
from sklearn.utils import check_random_state
from mnplib.miscoding import Miscoding

rng = check_random_state(42)

A = rng.randint(0, 2, size=6000)
B = rng.binomial(1, 0.1, size=6000)

X = np.column_stack((A, A, B))
y = 2 * A + B

names = np.array(["A", "A copy", "B"])

mi = partial(
    mutual_info_classif,
    discrete_features=True,
    random_state=42,
)

selector = SelectKBest(mi, k=2).fit(X, y)

print(
    "Mutual information:",
    names[selector.get_support()]
)

msd = Miscoding(
    X_type="categorical",
    y_type="categorical",
).fit(X, y)

selected = msd.select_features(max_features=2)

print("Miscoding:", names[selected])
\end{verbatim}}
\end{sourcecode}

A representative output is

\begin{sourcecode}
{\scriptsize \begin{verbatim}
Mutual information: ['A' 'A copy']
Miscoding:          ['A' 'B']
\end{verbatim}}
\end{sourcecode}

Individually, $A$ and its copy have exactly the same relation with the target, so a univariate mutual-information selector ranks both highly. Selecting both, however, provides no additional target information. Subset miscoding detects this because the second copy does not improve the joint representation. Adding $B$, on the other hand, completes the information required to reconstruct the target.
\end{example}

\subsection*{Miscoding of a Fitted Model}

The class can also evaluate the representation effectively used by a fitted machine-learning model \texttt{miscoding\_model(model)}. The library first identifies the subset of fitted features effectively used by
the model through its canonical model representation and then applies the same subset miscoding calculation described above. More detailed information can be obtained through \texttt{model\_analysis(model)}. This operation does not refit the model. It evaluates the representation selected by the fitted model relative to the dataset on which the \texttt{Miscoding} object was fitted. If the corresponding joint empirical distribution is unreliable, \texttt{miscoding\_model()} returns \texttt{NaN} and emits a warning rather
than returning a potentially misleading numerical value.

%
% Section: Inaccuracy
%

\section{Inaccuracy}
\label{sec:machine_learning_inaccuracy}

In Section~\ref{sec:inaccuracy:inaccuracy} we defined the inaccuracy of a description $d\in\mathcal{D}$ with respect to a representation $r\in\mathcal{R}$ as the normalized information distance between the intended representation $r$ and the representation $\Gamma(d)$ produced by the description:
\[
\iota(d,r) = \frac{\max\{K(r\mid\Gamma(d)), K(\Gamma(d)\mid r)\}}{\max\{K(r), K(\Gamma(d))\}}.
\]
Inaccuracy measures how much information is required to transform the output produced by a description into the intended representation, and conversely. A value close to zero means that the two representations can be transformed into each other using little additional information. A value close to one indicates that substantial information is required.

As with the other quantities based on Kolmogorov complexity, inaccuracy cannot be computed exactly in the general case. The \texttt{mnplib} library therefore replaces Kolmogorov complexity with the empirical code lengths introduced in Section~\ref{sec:utility_interface}.

In machine learning, the intended representation is the observed target vector $\mathbf{y} = (y_1,\ldots,y_n)$, while the output of the description is the vector of predictions produced by a fitted model, $\hat{\mathbf{y}} = (\hat y_1,\ldots,\hat y_n)$. The practical problem is therefore to measure how much information separates the observed targets from the predictions.

\subsection*{Empirical Inaccuracy}

Let $\hat K_C(\mathbf{y})$, $\hat K_C(\hat{\mathbf{y}})$, denote the empirical code lengths of the target and prediction vectors, and let $\hat K_C(\mathbf{y},\hat{\mathbf{y}})$ denote the empirical code length of their joint representation. As discussed in Section~\ref{sec:utility_interface}, these quantities are computed from the frequencies of the observed states. They do not correspond to an actual compressed bit stream. Using the usual approximation to normalized information distance, we obtain the practical inaccuracy used by \texttt{mnplib}.

Let $\mathbf{y}$ be a target vector and $\hat{\mathbf{y}}$ a vector of predictions of the same length. We define their empirical \emph{inaccuracy}\index{Inaccuracy} as
\[
\hat\iota(\mathbf{y},\hat{\mathbf{y}}) = \frac{ \hat K_C(\mathbf{y},\hat{\mathbf{y}})
- \min\{ \hat K_C(\mathbf{y}), \hat K_C(\hat{\mathbf{y}}) \} }{
\max\{ \hat K_C(\mathbf{y}), \hat K_C(\hat{\mathbf{y}}) \} }.
\]

The quantity is constrained in practice to the interval $[0,1]$. Lower values indicate that the predictions are informationally close to the observed target.

The expression is symmetric, $\hat\iota(\mathbf{y},\hat{\mathbf{y}}) = \hat\iota(\hat{\mathbf{y}},\mathbf{y})$, and therefore does not merely count how much information is required to correct the predictions into the target. It also considers how much information would be required in the converse direction. This distinguishes inaccuracy from conventional predictive metrics. Classification accuracy counts how many labels agree. Regression metrics such as mean squared error measure the magnitude of numerical residuals. Inaccuracy instead evaluates the information structure of the relationship between the complete target and prediction vectors.

The target determines how empirical code lengths are computed. If the target is categorical, both the target and predictions are treated as categorical attributes. Their distinct observed values are encoded as states,
their frequencies are counted, and the corresponding empirical code lengths are computed directly. If the target is numeric, the target and prediction vectors must first be discretized. The current implementation uses the automatic one-dimensional resolution introduced in Section~\ref{sec:utility_interface}.

The same resolved number of bins $b$ is used to compute $\hat K_C(\mathbf{y})$, $\hat K_C(\hat{\mathbf{y}})$, $\hat K_C(\mathbf{y},\hat{\mathbf{y}})$. This consistency is important because these quantities are subsequently
subtracted and divided. Using different resolutions would make part of the measured difference an artefact of discretization. 

When a prediction vector is already available, no feature matrix is required. The target can be fitted directly with

\begin{sourcecode}
{\scriptsize \begin{verbatim}
from mnplib.inaccuracy import Inaccuracy

metric = Inaccuracy().fit_y(y)

value = metric.inaccuracy_predictions(predictions)
\end{verbatim}}
\end{sourcecode}

The method \texttt{fit\_y()} stores the target, infers its type, computes its empirical code length, and prepares the estimator for comparison with one or more prediction vectors.

When a fitted machine-learning model is to be evaluated instead, the feature matrix is also required:

\begin{sourcecode}
{\scriptsize \begin{verbatim}
metric = Inaccuracy().fit(X, y)

value = metric.inaccuracy_model(model)
\end{verbatim}}
\end{sourcecode}

The model is not fitted again. It must provide a \texttt{predict()} method, which is called on the evaluation data stored by the \texttt{Inaccuracy} object.

A more detailed report is available through

\begin{sourcecode}
{\scriptsize \begin{verbatim}
metric.prediction_analysis(predictions)
metric.model_analysis(model)
\end{verbatim}}
\end{sourcecode}

\subsection*{The Structure of Prediction Errors}

Two prediction vectors may contain exactly the same number of incorrect predictions while differing substantially in the structure of those errors. This is an important difference between inaccuracy and the standar classification accuracy. Accuracy considers only whether each prediction is equal to the corresponding target. Inaccuracy considers the complete empirical relation between target and prediction states.

\begin{example}
\label{ex:machine_learning:inaccuracy:structured_errors}

Consider a balanced four-class target. We construct two predictors that misclassify exactly $20\%$ of the observations. In the first predictor, every error follows the same rule: $\hat y=(y+1)\bmod 4$. In the second predictor, each erroneous prediction is replaced by a randomly selected incorrect class.

\begin{sourcecode}
{\scriptsize \begin{verbatim}
import numpy as np
import pandas as pd

from sklearn.metrics import accuracy_score
from mnplib.inaccuracy import Inaccuracy

rng = np.random.default_rng(42)

n = 4000
y = np.tile(np.arange(4), n // 4)

wrong = np.zeros(n, dtype=bool)
wrong[::5] = True

systematic = y.copy()
systematic[wrong] = (y[wrong] + 1) % 4

random = y.copy()
for i in np.flatnonzero(wrong):
    random[i] = rng.choice(np.delete(np.arange(4), y[i]))

metric = Inaccuracy().fit_y(y)

print("Systematic bias",
      "Accuracy:", 1 - accuracy_score(y, systematic),
      "Inaccuracy:", metric.inaccuracy_predictions(systematic))

print("Random noise",
      "Accuracy:", 1 - accuracy_score(y, random),
      "Inaccuracy:", metric.inaccuracy_predictions(random))
\end{verbatim}}
\end{sourcecode}

A representative result is

\begin{center}
\begin{tabular}{lrr}
\hline
Predictions & Error rate & Inaccuracy \\
\hline
systematic & 0.20 & 0.36 \\
random     & 0.20 & 0.52 \\
\hline
\end{tabular}
\end{center}

Both predictors have exactly the same error rate. However, the systematic errors have a simpler empirical correction structure, and therefore produce a lower inaccuracy.
\end{example}

A second consequence of using empirical joint code lengths is that inaccuracy has a finer granularity than classification accuracy. For a dataset containing $n$ observations, classification accuracy depends only
on the number $c$ of correct predictions, $\operatorname{accuracy} = c/n$, since $c\in\{0,1,\ldots,n\}$, accuracy can assume at most $n+1$ distinct values. All prediction vectors containing the same number of correct labels therefore receive exactly the same accuracy, regardless of how their incorrect predictions are distributed among the possible classes. Inaccuracy is based instead on the empirical frequencies of the joint states $(y_i,\hat y_i)$. Consequently, different patterns of errors may produce different inaccuracy values even when the number of errors remains unchanged.

\begin{example}
\label{ex:machine_learning:inaccuracy:granularity}

Consider five balanced classes. We keep the classification error fixed at $20\%$, but progressively distribute the erroneous predictions among an increasing number of different wrong classes.

\begin{sourcecode}
{\scriptsize \begin{verbatim}
import numpy as np
import pandas as pd

from sklearn.metrics import accuracy_score
from mnplib.inaccuracy import Inaccuracy

n_classes        = 5
samples_per_class = 600
n_errors          = 120

y = np.repeat(np.arange(n_classes), samples_per_class)

metric = Inaccuracy().fit_y(y)

results = []

for n_error_types in range(1, n_classes):

    predictions = y.copy()

    for cls in range(n_classes):

        indices = np.flatnonzero(y == cls)[:n_errors]
        groups  = np.array_split(indices, n_error_types)

        for offset, group in enumerate(groups, start=1):
            predictions[group] = (cls + offset) % n_classes

    report = metric.prediction_analysis(predictions)

    results.append({
        "error_types": n_error_types,
        "accuracy": accuracy_score(y, predictions),
        "joint_states": report["n_observed_joint_states"],
        "inaccuracy": report["inaccuracy"],
    })

pd.DataFrame(results)
\end{verbatim}}
\end{sourcecode}

The result is approximately

\begin{center}
\begin{tabular}{rrrr}
\hline
Error types & Accuracy & Joint states & Inaccuracy \\
\hline
1 & 0.80 & 10 & 0.311 \\
2 & 0.80 & 15 & 0.397 \\
3 & 0.80 & 20 & 0.447 \\
4 & 0.80 & 25 & 0.483 \\
\hline
\end{tabular}
\end{center}

The number of incorrect predictions never changes, and therefore accuracy remains exactly $0.80$ throughout the experiment. What changes is the structure of the errors. With one error type, each true class is associated with only one additional incorrect prediction state. As more error types are introduced, the same errors are distributed among a larger number of target--prediction pairs. The joint empirical representation therefore requires a longer description, and inaccuracy increases.
\end{example}

Classical accuracy can also be misleading when the target distribution is strongly imbalanced. A model may achieve high accuracy simply by predicting the majority class while containing little information about the minority class.

\begin{example}
\label{ex:machine_learning:inaccuracy:unbalanced_dataset}

We create a binary classification problem in which $95\%$ of the observations belong to one class and $5\%$ to the other.

\begin{sourcecode}
{\scriptsize \begin{verbatim}
import numpy as np
import matplotlib.pyplot as plt

from sklearn.datasets import make_classification
from sklearn.tree import DecisionTreeClassifier
from mnplib.inaccuracy import Inaccuracy

X, y = make_classification(
    n_samples=1000,
    n_features=2,
    n_informative=2,
    n_redundant=0,
    class_sep=2,
    flip_y=0,
    weights=[0.95, 0.05],
    random_state=42,
)

metric = Inaccuracy().fit(X, y)

error = []
inacc = []

for size in np.arange(1, 100):

    model = DecisionTreeClassifier(
        min_samples_leaf=size,
        random_state=42,
    ).fit(X, y)

    error.append(1 - model.score(X, y))
    inacc.append(metric.inaccuracy_model(model))

plt.plot(np.arange(1, 100), error,
         label="1 - accuracy")
plt.plot(np.arange(1, 100), inacc,
         label="inaccuracy")
plt.xlabel("Minimum samples per leaf")
plt.ylabel("Error")
plt.legend()
plt.show()
\end{verbatim}}
\end{sourcecode}

As \texttt{min\_samples\_leaf} increases, the tree becomes less able to isolate the minority class. The ordinary classification error may remain close to the minority proportion because most observations still belong to the majority class. Inaccuracy, however, reflects the fact that the resulting prediction vector contains insufficient information to identify the minority observations.

\begin{figure}[h]
\centering
\includegraphics[width=0.6\textwidth]{inaccuracy_imbalanced.png}
\caption{Accuracy-based error and inaccuracy on an imbalanced dataset.}
\label{figure:machine_learning:inaccuracy:imbalanced}
\end{figure}

\end{example}

The same principle applies to regression. Metrics such as mean squared error and root mean squared error measure the numerical magnitude of residuals, but do not distinguish between residual patterns that have different informational structure.

\begin{example}
\label{ex:machine_learning:inaccuracy:same_rmse}

Consider two prediction vectors with exactly the same root mean squared error. The first has a constant systematic offset, $\hat y_i=y_i+1$, whereas the second adds irregular random noise having the same root mean
squared magnitude.

\begin{sourcecode}
{\scriptsize \begin{verbatim}
import numpy as np

from sklearn.metrics import mean_squared_error
from mnplib.inaccuracy import Inaccuracy

rng = np.random.default_rng(42)

n = 5000
y = rng.normal(0, 2, n)

biased = y + 1.0

noise  = rng.normal(size=n)
noise /= np.sqrt(np.mean(noise ** 2))
noisy  = y + noise

metric = Inaccuracy().fit_y(y)

print("Constant bias",
      "RMSE:", np.sqrt(mean_squared_error(y, biased)),
      "Inaccuracy:", metric.inaccuracy_predictions(biased))

print("Random noise",
      "RMSE:", np.sqrt(mean_squared_error(y, noisy)),
      "Inaccuracy:", metric.inaccuracy_predictions(noisy))
\end{verbatim}}
\end{sourcecode}

A representative result is

\begin{sourcecode}
{\scriptsize \begin{verbatim}
Constant bias RMSE: 1.0 Inaccuracy: 0.0
Random noise  RMSE: 1.0 Inaccuracy: 0.73
\end{verbatim}}
\end{sourcecode}

RMSE assigns the same value to both prediction vectors by construction. The constant offset, however, preserves the complete structure of the target. Under the equal-width discretization used by \texttt{mnplib}, translating all values by the same constant also translates the interval boundaries. The target and prediction vectors therefore receive the same sequence of discrete states and their empirical inaccuracy is zero. Random noise behaves differently. Different observations require different corrections, the target and prediction states no longer coincide, and the joint code length increases substantially.
\end{example}

This example illustrates a fundamental difference between the metrics. RMSE measures the magnitude of numerical errors. Inaccuracy measures how much information separates the prediction representation from the observed target.

%
% Section: Surfeit
%

\section{Surfeit}
\label{sec:machine_learning_surfeit}

In Chapter~\ref{chap:Surfeit} we introduced surfeit as a practical approximation to the surplus of a description. Let $d\in\mathcal{D}$ be a description. Its theoretical surfeit is
\[
\sigma(d) = 1 - \frac{K(d)}{l(d)},
\]
where $l(d)$ is the length of the description and $K(d)$ is its Kolmogorov complexity. Surfeit measures how much of the description can, in principle, be removed by a shorter description of the description itself. If $K(d)\approx l(d)$, then the description is close to incompressible and its surfeit is close to zero. If, on the other hand, $K(d)\ll l(d)$, then the description contains substantial compressible structure and its surfeit
is high.

Surfeit should not be confused with the theoretical surplus of a description. The latter is relative to the representation being described and asks how much information in a description is unnecessary with respect to a minimal description of that representation. Surfeit asks a weaker, but operationally useful, question: how much compressible structure remains in the description itself?

This distinction is particularly useful in machine learning. A trained model $m$ is represented by a finite description $d_m$. The model may predict the target accurately while still containing unnecessary branches, parameters, rules, repeated structures, or other descriptive machinery. Surfeit attempts to detect this remaining descriptive excess.

Unlike inaccuracy, surfeit does not evaluate the predictions produced by the model. Inaccuracy evaluates the relation between the model output and the observed target; surfeit evaluates the model description itself. Consequently, the two quantities are complementary. A very simple model may have low surfeit and high inaccuracy, while a highly elaborate model may have low inaccuracy and high surfeit.

% Practical Surfeit

\subsection*{Practical Surfeit}

Kolmogorov complexity cannot be computed in the general case. The \texttt{mnplib} library therefore approximates $K(d_m)$ by compressing a canonical string representation of the model. Let $L(d_m)$ denote the length, in bits, of the UTF-8 representation of the model description, and let $C(d_m)$ denote the length, also in bits, obtained after compressing that representation with \texttt{zlib}. The current implementation uses a fixed compression level of $9$. Keeping the compressor and its configuration fixed is important: practical surfeit is relative to the compression procedure, and changing the compressor may change the regularities that are detected.

A practical complication arises for short descriptions. The compressed string may be longer than the original one because a \texttt{zlib} stream contains a fixed wrapper overhead. Thus, we may observe $C(d_m)>L(d_m)$ even for a description containing some redundancy. The implementation assumes a fixed \texttt{zlib} overhead of six bytes, or $h_{\mathrm{zlib}}=48$ bits. It therefore defines the effective compressed length as
\[
C_{\mathrm{eff}}(d_m) = \min\left\{L(d_m), \max\left\{0, C(d_m)-h_{\mathrm{zlib}}\right\}\right\}.
\]
The first clipping operation prevents the corrected compressed length from becoming negative. The second ensures that compression is never treated as making the description longer than its original representation.

The \texttt{mnplib} implementation introduces an additional target-dependent safeguard. Let $\hat K_C(\mathbf{y})$ be the empirical code length of the target attribute, computed using the procedures introduced in Section~\ref{sec:utility_interface}. The reference description length is defined as
\[
L_{\mathrm{ref}}(m,\mathbf{y}) = \min\left\{C_{\mathrm{eff}}(d_m), \hat K_C(\mathbf{y}) \right\}.
\]
This condition prevents the compressed model description from being credited with more useful descriptive information than is present in the target that the model is intended to describe.

The theoretical surfeit $\sigma(d)$ depends only on the description. The quantity implemented by \texttt{mnplib} is therefore a target-calibrated practical approximation: the internal compressibility of the model remains the primary quantity, but the useful reference length is bounded by the empirical information contained in the target.

Let $m$ be a fitted model, let $d_m$ be its canonical model description, and let $\mathbf{y}$ be the target attribute. We define the empirical \emph{surfeit}\index{Surfeit} of $m$ with respect to $\mathbf{y}$ as
\[
\hat\sigma(m,\mathbf{y}) = 1 - \frac{L_{\mathrm{ref}}(m,\mathbf{y})}{L(d_m)},
\]
where
\[
L_{\mathrm{ref}}(m,\mathbf{y}) = \min\left\{C_{\mathrm{eff}}(d_m), \hat K_C(\mathbf{y})\right\}.
\]
The resulting value belongs to the interval $[0,1]$. Low values indicate that the model description contains little compressible excess relative to the chosen coding and compression procedures. High values indicate that a substantial fraction of the description is either compressible or excessive relative to the information contained in the target. For categorical targets, $\hat K_C(\mathbf{y})$ is computed directly from the frequencies of the observed target states. Numeric targets are discretized using the automatic one-dimensional resolution introduced in Section~\ref{sec:utility_interface}.

The most direct use of the Surfeit class is to evaluate a fitted estimator:

\begin{sourcecode}
{\scriptsize
\begin{verbatim}
from mnplib.surfeit import Surfeit

sft = Surfeit().fit(X, y)

surfeit = sft.surfeit_model(model)
\end{verbatim}
}
\end{sourcecode}

The method \texttt{surfeit\_model()} does not serialize the Python object using \texttt{str()}, \texttt{repr()}, or a binary object serializer. Instead, the estimator is passed through the model adapter layer, which produces its canonical description string. The resulting string is then evaluated using exactly the same procedure as any explicit description.

When the model description is already available, the feature matrix is not required:

\begin{sourcecode}
{\scriptsize
\begin{verbatim}
from mnplib.surfeit import Surfeit

sft = Surfeit().fit_y(y)

surfeit = sft.surfeit_string(model_string)
\end{verbatim}
}
\end{sourcecode}

Thus, \texttt{surfeit\_string()} is the primitive metric operation, while \texttt{surfeit\_model()} adds the model-serialization step.

A detailed explanation of an explicit description is available through

\begin{sourcecode}
{\scriptsize
\begin{verbatim}
details = sft.description_analysis(model_string)
\end{verbatim}
}
\end{sourcecode}

The returned dictionary contains

\begin{sourcecode}
{\scriptsize
\begin{verbatim}
surfeit
model_code_length_bits
compressed_code_length_bits
effective_compressed_code_length_bits
target_code_length_bits
reference_code_length_bits
compression_ratio
reference_source
\end{verbatim}
}
\end{sourcecode}

The field \texttt{compression\_ratio} is the ratio between the raw \texttt{zlib}-compressed size and the raw model size before overhead correction. It may therefore be greater than one for short descriptions. The field \texttt{reference\_source} indicates which quantity limits the reference length. Its value is \texttt{"compression"} when the corrected compressed description is shorter than the target code length, \texttt{"target"} when the target code length is shorter, and \texttt{"both"} when they coincide.

For a fitted estimator,

\begin{sourcecode}
{\scriptsize
\begin{verbatim}
details = sft.model_analysis(model)
\end{verbatim}
}
\end{sourcecode}

returns the same surfeit diagnostics together with the canonical model string, the model type, the selected feature indices, and the number of selected features.

A functional interface is also provided:

\begin{sourcecode}
{\scriptsize
\begin{verbatim}
from mnplib.surfeit import surfeit_model

value = surfeit_model(model, X=X, y=y)
\end{verbatim}
}
\end{sourcecode}

% A Minimal Language for Model Descriptions

\subsection*{A Minimal Language for Model Descriptions}

Computing surfeit requires converting a trained model into a string. This step is more important than it may initially appear. Suppose that the same model is encoded using two different languages. One language may provide compact operators for a particular model family, while another may require a verbose representation of the same computation. The resulting description lengths and compression ratios could differ even though the underlying predictive procedure is identical. This problem becomes particularly important when comparing models from different families. A decision tree, a linear model, and a neural network must not be compared using unrelated serialization formats. Otherwise, the measured difference in surfeit might reflect the serialization conventions rather than the models themselves.

For this reason, \texttt{mnplib} uses a model adapter layer that converts supported fitted estimators into canonical textual descriptions. We assume that these descriptions are expressed in a common restricted Python-like language. A proposed grammar for this language is

\begin{sourcecode}
{\scriptsize
\begin{verbatim}
<program>   ::= <function> | <function> <program>
<function>  ::= "def" NAME "(" <params> ")" ":" <suite>
<params>    ::= epsilon | NAME | NAME "," <params>
<suite>     ::= NEWLINE INDENT <block> DEDENT
<block>     ::= <statement> | <statement> <block>

<statement> ::= <ref> "=" <expr> NEWLINE
              | <ref> ".append(" <expr> ")" NEWLINE
              | "return" <expr> NEWLINE
              | "if" <test> ":" <suite> <else>
              | "for" NAME "in" "range(" <expr> "):" <suite>

<else>      ::= epsilon | "else:" <suite>

<expr>      ::= NUMBER | <ref> | "[" <args> "]" | "-" <expr>
              | "(" <expr> <op> <expr> ")"
              | NAME "(" <args> ")"

<args>      ::= epsilon | <expr> | <expr> "," <args>
<ref>       ::= NAME | <ref> "[" <expr> "]"
<test>      ::= <expr> <comparison> <expr>

<op>        ::= "+" | "-" | "*" | "/"
<comparison> ::= "<" | "<=" | ">" | ">=" | "==" | "!="
\end{verbatim}
}
\end{sourcecode}

The language contains only the constructs required to express predictive procedures: functions, assignments, arrays, arithmetic expressions, function calls, conditionals, loops, indexing, and return statements. It deliberately excludes implementation details that are irrelevant to the predictive procedure, such as library-specific object metadata, memory addresses, caches, training history, or debugging information.

The grammar alone, however, is not sufficient to guarantee comparable model descriptions. The serializer should also use deterministic conventions for variable names, numerical precision, statement ordering, indentation, and whitespace. Otherwise, two equivalent implementations could receive different surfeit values merely because of superficial syntactic differences.

The resulting surfeit should therefore always be interpreted relative to this common description language. This is analogous to fixing a reference universal Turing machine when discussing Kolmogorov complexity: the exact numerical value depends on the coding convention, but comparisons become meaningful once that convention is held fixed.

\subsection*{Models from the Same Family}

Surfeit can distinguish between models from the same family that obtain similar predictive performance but have substantially different descriptive complexity.

\begin{example}
\label{ex:machine_learning:surfeit:same_family}

Consider the Breast Cancer Wisconsin dataset included in \texttt{scikit-learn}. We train two decision trees, one restricted to a maximum depth of two and another allowed to grow to depth five.

\begin{sourcecode}
{\scriptsize
\begin{verbatim}
from sklearn.datasets import load_breast_cancer
from sklearn.model_selection import train_test_split
from sklearn.tree import DecisionTreeClassifier

from mnplib.surfeit import surfeit_model

X, y = load_breast_cancer(return_X_y=True)

X_train, X_test, y_train, y_test = train_test_split(
    X, y, test_size=0.2, random_state=42
)

simple = DecisionTreeClassifier(
    max_depth=2,
    random_state=42,
).fit(X_train, y_train)

complex_ = DecisionTreeClassifier(
    max_depth=5,
    random_state=42,
).fit(X_train, y_train)

print("Simple accuracy:", simple.score(X_test, y_test))
print("Complex accuracy:", complex_.score(X_test, y_test))

print(
    "Simple surfeit:",
    surfeit_model(simple, X=X_train, y=y_train),
)
print(
    "Complex surfeit:",
    surfeit_model(complex_, X=X_train, y=y_train),
)
\end{verbatim}}
\end{sourcecode}

A representative result is

\begin{sourcecode}
{\scriptsize
\begin{verbatim}
Simple accuracy:  0.930
Complex accuracy: 0.947

Simple surfeit:   0.560
Complex surfeit:  0.895
\end{verbatim}}
\end{sourcecode}

The deeper tree improves the test accuracy by less than two percentage points, but its surfeit is substantially larger. Conventional predictive performance alone does not expose this difference in descriptive economy.
\end{example}

This example also illustrates why surfeit is not simply the number of parameters, tree nodes, or model depth. Those quantities are useful within a particular model family, but they are not directly comparable across different kinds of models. Surfeit evaluates the resulting descriptions through the same coding principle.

% The Inaccuracy--Surfeit Trade-off

\subsection*{The Inaccuracy--Surfeit Trade-off}

Increasing model complexity often decreases inaccuracy because a more flexible model can reproduce a larger fraction of the structure present in the target. At the same time, the resulting description generally becomes more elaborate. The two effects can therefore move in opposite directions.

\begin{example}
\label{ex:machine_learning:surfeit:polynomial}

We generate $900$ observations from the sinusoidal relation $y=\cos(1.5\pi x)$ and fit polynomial models of increasing degree. Each polynomial model is implemented as a linear regression over polynomial features.

\begin{sourcecode}
{\scriptsize
\begin{verbatim}
import numpy as np
import pandas as pd
import matplotlib.pyplot as plt

from sklearn.linear_model import LinearRegression
from sklearn.preprocessing import PolynomialFeatures

from mnplib.inaccuracy import Inaccuracy
from mnplib.surfeit import Surfeit

rng = np.random.default_rng(42)

X = np.sort(
    rng.uniform(0.0, 3.0, size=900)
).reshape(-1, 1)

y = np.cos(1.5 * np.pi * X[:, 0])

inacc = Inaccuracy().fit_y(y)

rows = []

for degree in range(1, 8):

    poly = PolynomialFeatures(
        degree=degree,
        include_bias=False,
    )
    X_poly = poly.fit_transform(X)

    model = LinearRegression().fit(X_poly, y)
    predictions = model.predict(X_poly)

    sft = Surfeit().fit(X_poly, y)

    rows.append({
        "degree": degree,
        "inaccuracy":
            inacc.inaccuracy_predictions(predictions),
        "surfeit":
            sft.surfeit_model(model),
    })

results = pd.DataFrame(rows)

plt.plot(
    results["degree"],
    results["inaccuracy"],
    label="Inaccuracy",
)

plt.plot(
    results["degree"],
    results["surfeit"],
    label="Surfeit",
)

plt.xlabel("Polynomial degree")
plt.ylabel("Metric")
plt.legend()
plt.show()
\end{verbatim}}
\end{sourcecode}

\begin{figure}[ht]
\centering
\includegraphics[width=0.6\textwidth]{surfeit_vs_inaccuracy.png}
\caption{Inaccuracy and surfeit for polynomial models of increasing degree.}
\label{figure:machine_learning:surfeit:tradeoff}
\end{figure}

Low-degree polynomials are too restrictive to reproduce the sinusoidal relation and consequently have relatively high inaccuracy. As the degree increases, inaccuracy decreases because the polynomial becomes more expressive.

At the same time, the fitted description becomes increasingly elaborate and surfeit tends to increase. Beyond some point, additional complexity may produce only a small reduction in inaccuracy while introducing a substantial increase in surfeit.
\end{example}

The example illustrates the role of the two quantities. Inaccuracy encourages descriptions capable of reconstructing the observed target. Surfeit encourages descriptions that do so economically. Neither metric should be minimized in isolation; their trade-off is one of the problems addressed by the nescience criterion introduced in the following section.

% Subsection: Comparing Different Model Families

\subsection*{Comparing Different Model Families}

A particularly useful property of surfeit is that it provides a common description-based criterion for models whose conventional complexity measures are not comparable. The depth of a decision tree cannot be directly compared with the number of weights in a neural network, nor with the number of coefficients in a linear model. Once the models are expressed in the same canonical language, however, their descriptions can be analyzed using the same surfeit calculation.

\begin{example}
\label{ex:machine_learning:surfeit:model_families}

Consider a binary classification problem consisting of two approximately separable Gaussian clusters. We train a logistic regression model, a decision tree, and a neural network.

\begin{sourcecode}
{\scriptsize
\begin{verbatim}
from sklearn.datasets import make_blobs
from sklearn.linear_model import LogisticRegression
from sklearn.model_selection import train_test_split
from sklearn.neural_network import MLPClassifier
from sklearn.tree import DecisionTreeClassifier

from mnplib.surfeit import Surfeit

X, y = make_blobs(
    n_samples=2000,
    centers=2,
    cluster_std=2,
    n_features=2,
    random_state=42,
)

X_train, X_test, y_train, y_test = train_test_split(
    X,
    y,
    test_size=0.5,
    stratify=y,
    random_state=42,
)

linear = LogisticRegression().fit(X_train, y_train)

tree = DecisionTreeClassifier(
    random_state=42,
).fit(X_train, y_train)

neural = MLPClassifier(
    hidden_layer_sizes=(20,),
    max_iter=2000,
    random_state=42,
).fit(X_train, y_train)

sft = Surfeit().fit(X_train, y_train)

print(
    "Linear:",
    linear.score(X_test, y_test),
    sft.surfeit_model(linear),
)

print(
    "Decision tree:",
    tree.score(X_test, y_test),
    sft.surfeit_model(tree),
)

print(
    "Neural network:",
    neural.score(X_test, y_test),
    sft.surfeit_model(neural),
)
\end{verbatim}}
\end{sourcecode}

A representative result is

\begin{center}
\begin{tabular}{lrr}
\hline
Model & Accuracy & Surfeit \\
\hline
Linear model   & 0.993 & 0.163 \\
Decision tree  & 0.991 & 0.636 \\
Neural network & 0.986 & 0.889 \\
\hline
\end{tabular}
\end{center}

All three models solve the classification problem with very similar predictive performance. Their descriptions, however, differ substantially.

For this simple decision boundary, the linear model provides the most economical description under the common model language. The decision tree requires a more elaborate set of rules, while the neural network requires a substantially larger collection of parameters and operations.
\end{example}


The conclusion is specific to the fitted models and the problem under study. Surfeit does not impose a general preference for linear models over decision trees or neural networks. On a sufficiently complex problem, a neural network may provide a more economical description than a very large tree or an elaborate collection of linear pieces.

The validity of this comparison depends critically on the use of the same canonical description language. If each model family were serialized according to unrelated conventions, the resulting surfeit values would partly measure the verbosity of those conventions rather than the descriptive complexity of the models themselves.

Surfeit is related to model complexity, and excessive model complexity is often associated with overfitting. Nevertheless, a high surfeit should not be interpreted as a statistical test for overfitting.

Surfeit analyzes the compressibility of the model description. It does not directly measure generalization error, and it does not compare training and test performance. A model may have high surfeit and nevertheless generalize well if the phenomenon genuinely requires a complex description. Conversely, a model may have low surfeit simply because it is too simple to describe the phenomenon.

What surfeit can identify is a situation in which increasing descriptive complexity appears difficult to justify by the corresponding reduction in inaccuracy. Such a pattern is often associated with overfitting, but the interpretation must be made jointly with the predictive behavior of the model.

%
% Section: Nescience
%

\section{Nescience}
\label{sec:machine_learning_nescience}

In Chapter \ref{chap:Nescience} we defined the concept of nescience as the solution to a nonlinear multiobjective optimization problem, in which we minimize the miscoding, inaccuracy, and surfeit of representations and models. The solution to this problem is, in general, not unique: we can often find multiple pairs (representation, model) such that none of the three quantities can be improved without degrading at least one of the others (Pareto optimality). However, in practice we expect a machine-learning library to provide a single solution when training a model on a dataset. To provide a unique solution, we resort to a utility function that selects one element from the Pareto-optimal set. The \texttt{nescience} library offers several utility functions; the default is the arithmetic mean of the three metrics.

\begin{remark}
The \texttt{nescience} library implements the following utility functions to approximate the concept of nescience, i.e., to compute $\hat\nu\!\left(\mathbb{Z}, m, \mathbf{y} \right)$:

\begin{itemize}
\item Euclidean distance: $\left( \hat\mu(\mathbb{Z}, \mathbf{y})^2 + \hat\iota(\hat{y}, \mathbf{y})^2  + \hat\sigma(m, \mathbf{y})^2 \right)^{1/2}$
\item Arithmetic mean: $\frac{\hat\mu(\mathbb{Z}, \mathbf{y}) + \hat\iota(\hat{y}, \mathbf{y}) + \hat\sigma(m, \mathbf{y})}{3}$
\item Geometric mean: $\left( \hat\mu(\mathbb{Z}, \mathbf{y}) \times \hat\iota(\hat{y}, \mathbf{y}) \times \hat\sigma(m, \mathbf{y}) \right)^{1/3}$
\item Product: $\hat\mu(\mathbb{Z}, \mathbf{y}) \times \hat\iota(\hat{y}, \mathbf{y}) \times \hat\sigma(m, \mathbf{y})$
\item Addition: $\hat\mu(\mathbb{Z}, \mathbf{y}) + \hat\iota(\hat{y}, \mathbf{y}) + \hat\sigma(m, \mathbf{y})$
\item Weighted mean: $w_\mu \hat\mu(\mathbb{Z}, \mathbf{y}) + w_\iota \hat\iota(\hat{y}, \mathbf{y}) + w_\sigma \hat\sigma(m, \mathbf{y})$
\item Harmonic mean: $\frac{3}{ \hat\mu(\mathbb{Z}, \mathbf{y})^{-1} + \hat\iota(\hat{y}, \mathbf{y})^{-1} + \hat\sigma(m, \mathbf{y})^{-1} }$
\end{itemize}

Euclidean distance and addition can produce nescience values greater than one, which conflicts with the intended $[0,1]$ range. The geometric mean, product, and harmonic mean have the issue that the resulting nescience is zero (or undefined, in the harmonic case) whenever any of the three components is zero. The weighted mean introduces three additional hyperparameters $(w_\mu, w_\iota, w_\sigma)$ that must be chosen (often subject to $w_\mu + w_\iota + w_\sigma = 1$). It remains an open question which utility function is most appropriate in general; by default we use the arithmetic mean.
\end{remark}

Example \ref{ex:nescience_decisiontree} shows how to use the \texttt{nescience} library to compute the nescience of a dataset and a model.

\begin{example}
\label{ex:nescience_decisiontree}

We compute the nescience of a decision-tree classifier on the \texttt{digits} dataset (MNIST handwritten digits) from \texttt{sklearn}.

\begin{sourcecode}
{\scriptsize \begin{verbatim}
from sklearn.tree import DecisionTreeClassifier
from sklearn.datasets import load_digits
from Nescience.Nescience import Nescience

data = load_digits()

tree = DecisionTreeClassifier()
tree.fit(data.data, data.target)
tree.score(data.data, data.target
[ ] 1

nescience = Nescience()
nescience.fit(data.data, data.target)

nescience.nescience(tree)
[ ] 0.5895603819965907
\end{verbatim}}
\end{sourcecode}

The training accuracy of the decision tree is $1.0$, meaning all samples are correctly classified on the training set. This is symptomatic of overfitting. In practice we would use a train/test split or cross-validation. Nonetheless, the nescience value of about $0.59$ flags that something is wrong with the model-data pairing.

\end{example}

As Example \ref{ex:nescience_decisiontree} illustrates, a key advantage of nescience is that it can evaluate model quality without requiring computationally expensive procedures such as cross-validation or holding out a test set. Another advantage is that nescience allows comparisons across different model families, as shown next.

\begin{example}
\label{ex:nescience_comparison}

We compare two model families on the Breast Cancer Wisconsin (Diagnostic) dataset from \texttt{sklearn}: a decision tree and a neural network.

\begin{sourcecode}
{\scriptsize \begin{verbatim}
from sklearn.tree import DecisionTreeClassifier
from sklearn.neural_network import MLPClassifier
from sklearn.datasets import load_breast_cancer
from Nescience.Nescience import Nescience

data = load_breast_cancer()
X = data.data
y = data.target

tree = DecisionTreeClassifier(max_depth=3)
tree.fit(X, y)
tree.score(X, y)
[ ] 0.9789103690685413

nescience = Nescience()
nescience.fit(X, y)
nescience.nescience(tree)
[ ] 0.5945936419010083

nn = MLPClassifier()
nn.fit(X, y)
nn.score(X, y)
[ ] 0.9261862917398945

nescience.nescience(nn)
[ ] 0.7860523786210711
\end{verbatim}}
\end{sourcecode}

In one run, both models achieve high and broadly similar training accuracy. In this case, not only does the decision tree achieve the better score, but its nescience is also much lower than that of the multilayer perceptron; hence we should prefer the former to the latter.

\end{example}

Nescience can also be used to optimize hyperparameters within a (parametric) family of models. A practical advantage is that we can adopt a greedy search for a hyperparameter that is monotone with accuracy: select the value at which nescience stops decreasing and starts to increase, since this is the point at which we are no longer learning anything new from the dataset (see Example \ref{ex:nescence_hyperparameter}).

\begin{example}
\label{ex:nescence_hyperparameter}

We revisit a decision tree on the breast cancer dataset. We train $10$ trees with \texttt{max\_depth} from $1$ to $10$. The deeper the tree, the higher the training accuracy, but also the higher the risk of overfitting. For each tree we compute the model's nescience and compare it to a cross-validation score. Figure \ref{figure:nescience_cancer} shows typical behavior: both nescience and the cross-validation error decrease as depth increases, up to an inflection point after which they increase, indicating overfitting. The \texttt{nescience} library suggests a maximum depth of $7$, whereas cross-validation indicates an optimum near $6$.

\begin{figure}[h]
\centering
\includegraphics[width=0.6\textwidth]{nescience_cancer.png}
\caption{Evolution of Nescience with Tree Depth.}
\label{figure:nescience_cancer}
\end{figure}

\end{example}

It is instructive to examine the three components in Figure \ref{figure:nescience_cancer}. As expected, deeper trees reduce inaccuracy and increase surfeit. By contrast, miscoding can evolve nonmonotonically because each candidate tree may use a different subset of features at internal nodes. It would be desirable to have a tree-building algorithm that accounts for miscoding when selecting split features; such an algorithm is described in Section \ref{sec:decision_trees}.

Finally, we consider hyperparameter searches where a greedy approach is not applicable—for example, when searching over multiple (often conflicting) hyperparameters. Grid search can be computationally expensive due to the number of combinations, and even more so if each candidate requires cross-validation. Since nescience can flag overfitting without cross-validation, it can substantially speed up hyperparameter search. Example \ref{ex:hyper_search} shows how to integrate nescience into \texttt{GridSearchCV}.

\begin{example}
\label{ex:hyper_search}

In this example we are going to see how we can use the \texttt{nescience} library to find the optimal hyperparameters for a model using a grid search. In particular, we are going to select the best hyperparameters for a multilayer perceptron classifer, including the number of hidden layers, and the size of those layers (what it is called Neural Architecture Search). The procedure will be demonstrated using the \texttt{digits} dataset.

\begin{sourcecode}
{\scriptsize \begin{verbatim}
from Nescience.Nescience import Nescience
from sklearn.neural_network import MLPClassifier
from sklearn.model_selection import GridSearchCV
from sklearn.metrics import classification_report
from sklearn.model_selection import train_test_split
\end{verbatim}}
\end{sourcecode}

First of all we have to provide a custom loss function based on the concept of nescience to be integrated with the search procedure. The next code shows how to implement such a function.

\begin{sourcecode}
{\scriptsize \begin{verbatim}
def my_custom_loss_func(estimator, X, y):
    
    nsc = Nescience()
    nsc.fit(X, y)
    nescience = nsc.nescience(estimator)
    
    # scikit-learn expect that higher numbers are better
    score = -nescience
    
    return score
\end{verbatim}}
\end{sourcecode}

Second, we have to define the grid of hyperparameters over which we are going to do the search. The larger the grid, the better the result, but also, the more computer time is required to evaluate all possible combinations.

\begin{sourcecode}
{\scriptsize \begin{verbatim}
parameters = {'solver': ['lbfgs'],
              'max_iter': [1000, 1500, 2000 ], 
              'alpha': 10.0 ** -np.arange(1, 10, 3),
              'hidden_layer_sizes':[(60,), (100,), (60, 60,), (100, 100,), 
                                    (60, 60, 60,), (100, 100, 100,)]}
\end{verbatim}}
\end{sourcecode}

Next code show how to do a classical grid search using the score of the models. The search will be evaluated using a train/test split of the dataset.
 
\begin{sourcecode}
{\scriptsize \begin{verbatim}
clf_std = GridSearchCV(estimator=MLPClassifier(), param_grid=parameters,
                       cv=3, iid=True, n_jobs=-1)
clf_std.fit(X_train, y_train)
clf_std.best_params_

[] {'alpha': 0.1,
[]  'hidden_layer_sizes': (100,),
[]  'max_iter': 1000,
[]  'solver': 'lbfgs'}

y_true, y_pred = y_test, clf_std.predict(X_test)
print(classification_report(y_true, y_pred))

[] precision    recall  f1-score   support
[] avg / total       0.98      0.97      0.97       540
\end{verbatim}}
\end{sourcecode}

Next code show how to perform exactly the same search, but using the concept of nescience instead of the metric score.

\begin{sourcecode}
{\scriptsize \begin{verbatim}
clf_nsc = GridSearchCV(estimator=MLPClassifier(), param_grid=parameters,
                       cv=3, scoring=my_custom_loss_func, iid=True)
clf_nsc.fit(X_train, y_train)
clf_nsc.best_params_

{'alpha': 0.1,
 'hidden_layer_sizes': (60,),
 'max_iter': 1500,
 'solver': 'lbfgs'}

y_true, y_pred = y_test, clf_nsc.predict(X_test)
print(classification_report(y_true, y_pred))

                  precision    recall    f1-score    support
avg / total       0.98         0.98      0.98        540
\end{verbatim}}
\end{sourcecode}

In a typical run, both searches yield strong generalization. Notably, the nescience-based search selects a \emph{smaller} model (e.g., a single hidden layer with $60$ neurons instead of $100$) and compensates with a higher \texttt{max\_iter}. This reflects the role of surfeit: nescience tends to select the smallest model that attains high accuracy without overfitting the training data.

\end{example}

%
% Section: Auto Classification
%

\section{Auto Classification}
\label{sec:machine_learning_classification}

The \texttt{nescience} library also includes a module for automated machine learning (for both classification and regression problems). The AutoML module returns the single model, chosen from a collection of model families, that yields the smallest nescience. For each family, the class performs a greedy search over the hyperparameters required by that family, using nescience as the objective.

The next example shows how to apply the AutoML tools.

\begin{example}
\label{ex:automl}

In this example we use the \texttt{nescience} library to find the model that best describes the \texttt{digits} dataset.

\begin{sourcecode}
{\scriptsize \begin{verbatim}
from sklearn.datasets import load_digits
from sklearn.model_selection import train_test_split

from Nescience.Nescience import AutoClassifier

(X, y) = load_digits(return_X_y=True)
X_train, X_test, y_train, y_test = train_test_split(X, y, random_state=1)

model = AutoClassifier()
model.fit(X_train, y_train)

model.score(X_test, y_test)
[] 0.9622222222222222
\end{verbatim}}

\end{sourcecode}
\end{example}

If we evaluate \texttt{type(model.model)}, we will see that the library has selected a linear support vector machine as the best model for this dataset.

A key difference between the \texttt{nescience} library and other AutoML libraries is that it returns a \emph{single} model as the best candidate, rather than an ensemble of models. In this way, the data scientist can directly reuse the result of \texttt{AutoClassifier} and proceed with the subsequent analysis.

%
% Section: Auto Regression
%

\section{Auto Regression}

The \texttt{AutoRegressor} class automatically selects the best model for a regression problem. In particular, it (i) computes an optimal subset of features, (ii) selects the most suitable family of models, and (iii) tunes the model's hyperparameters—using nescience as the objective. 

\begin{example}
\label{ex:auto_regression}

In this example we apply our auto-regression class to estimating house prices, using the \texttt{boston} dataset included with \texttt{scikit-learn}.

\begin{sourcecode}
{\scriptsize \begin{verbatim}
from fastautoml.fastautoml import AutoRegressor
from sklearn.datasets import load_boston
from sklearn.model_selection import train_test_split

(X, y) = load_boston(return_X_y=True)
X_train, X_test, y_train, y_test = train_test_split(X, y, random_state=1)

model = AutoRegressor()
model.fit(X_train, y_train)
AutoRegressor()

model.score(X_test, y_test)
0.8763987309111113
\end{verbatim}}

\end{sourcecode}
\end{example}

If we evaluate \texttt{type(model.model)}, we see that the library has selected a linear support vector machine or a decision tree (depending on the run and dataset version) as the best model for this dataset.

A key difference between the \texttt{nescience} library and other AutoML libraries is that it returns a single model as the best candidate, rather than an ensemble of models. In this way, the data scientist can directly reuse the result of \texttt{AutoRegressor} and proceed with the analysis.

%
% Section: Time Series
%

\section{Time Series}

In this section we study the application of the concept of miscoding to a time series and a delayed version of itself, or to a delayed version of a second time series.

Automiscoding applies miscoding to a time series and a lagged version of itself, as a function of the lag. Automiscoding estimates to what \emph{extent} past observations of the series can explain (or help forecast) future observations. In this sense, automiscoding has a similar objective to autocorrelation in classical time-series analysis (see Section \ref{sub:autocorrelation}).

\begin{definition}
Let $\{\mathbf{x}_t\}$ be a time series composed by $n$ samples. We define lag $k$ \emph{regular automiscoding} of $\{\mathbf{x}_t\}$ as $\hat\mu(\mathbf{x}_{x_{k+1}, x_{k+2}, \ldots, x_n}, \mathbf{x}_{x_0, x_1, \ldots, x_{(n-k)}})$. We define in the same way the concepts of \emph{adjusted automiscoding} and \emph{partial automiscoding}.
\end{definition}

In contrast to autocorrelation, automiscoding is defined for all time series, including those with a trend. Moreover, automiscoding remains interpretable in the presence of trends and can reveal seasonal components without requiring prior decomposition.

\begin{example}
We examine cycles in the number of passengers of a US airline. Figure \ref{figure:air_passengers} shows monthly passengers from 1949 to 1960 (AirPassengers dataset, see References below). There is a clear yearly cycle. We apply automiscoding to confirm this analytically.

\begin{figure}[h]
\centering
\includegraphics[width=0.6\textwidth]{airpassengers.png}
\caption{Air Passengers.}
\label{figure:air_passengers}
\end{figure}

\begin{sourcecode}
{\scriptsize \begin{verbatim}
from nescience.timeseries import TimeSeries

data = pd.read_csv("data/AirPassengers.csv", index_col=["Month"],
                    parse_dates=True)
X = np.array(data["#Passengers"]).reshape(-1, 1)

ts = TimeSeries(auto=False)
ts.fit(data)
mscd = ts.auto_miscoding(max_lag=36)
\end{verbatim}}
\end{sourcecode}

As shown in Figure \ref{figure:auto-miscoding}, the adjusted automiscoding exhibits a prominent peak every twelve months: the distance between the aligned sequences is minimized when the lag is a multiple of one year.

\begin{figure}[h]
\centering
\includegraphics[width=0.6\textwidth]{auto-miscoding.png}
\caption{Auto-miscoding of Air Passengers.}
\label{figure:auto-miscoding}
\end{figure}

\end{example}

Crossmiscoding measures the relationship between a time series and a lagged version of a second time series. The objective is to detect whether the first series has temporal predictive power over the second.

\begin{definition}
Let $\{\mathbf{x}_t\}$ and $\{\mathbf{y}_t\}$ be two time series composed by $n$ samples each. We define lag $k$ \emph{regular crossmiscoding} of $\{\mathbf{x}_t\}$ and $\{\mathbf{y}_t\}$ as $\hat\mu(\mathbf{x}_{x_{k+1}, x_{k+2}, \ldots, x_n}, \mathbf{y}_{x_0, x_1, \ldots, x_{(n-k)}})$. We define in the same way the concepts of \emph{adjusted crossmiscoding} and \emph{partial crossmiscoding}.
\end{definition}

\begin{example}
We investigate whether it is possible to predict household appliance energy consumption. The dataset (Appliances Energy Prediction; see References) contains temperature and humidity measurements from rooms every ten minutes, the appliance energy consumption, and weather variables from a nearby station.

For each feature we compute the lag in $[1,30]$ that maximizes crossmiscoding with the target:

\begin{sourcecode}
{\scriptsize \begin{verbatim}
from fastautoml.fastautoml import Miscoding

X = pd.read_csv("../data/energydata_complete.csv", parse_dates=["date"],
                index_col="date")
y = X["Appliances"]
X = X.drop(["Appliances", "lights"], axis=1)

miscoding = Miscoding()
miscoding.fit(X, y)

best_lag  = list()
for i in np.arange(X.shape[1]):
    mscd = miscoding.cross_miscoding(attribute1=i, min_lag=1, max_lag=30)
    best_lag.append(np.where(mscd == np.max(mscd))[0][0] + 1)
\end{verbatim}}
\end{sourcecode}

Figure \ref{figure:cross-miscoding} shows the results. In general, indoor measurements favor small lags, whereas weather variables (which influence indoor conditions with delay) favor larger lags.

\begin{figure}[h]
\centering
\includegraphics[width=0.6\textwidth]{cross_miscoding_lag.png}
\caption{Cross Miscoding Lag}
\label{figure:cross-miscoding}
\end{figure}

\end{example}

\begin{remark}
The approximation to miscoding introduced in this chapter estimates (i) the quality of individual features as predictors of a target and (ii) the quality of the training dataset as a whole. However, it does not explicitly account for redundancy among the features themselves. For instance, two features $\mathbf{x}_i$ and $\mathbf{x}_j$ may each have low miscoding with respect to $\mathbf{y}$ yet be largely redundant with one another. It remains an open question how to extend miscoding to incorporate feature redundancy in a way that stays close to the theoretical definition, is computationally efficient, and does not require prohibitively large samples.
\end{remark}

% Auto Time Series

% \subsection{Auto Time Series}

% {\color{red} TODO: Introduce the auto-time series models, and clearly state what it can be expect from such models (short story: nothing)}

% {\color{red} TODO: Introduce structural time series models, the state space representation, and the Kalman filter}

% {\color{red} structura approach [...] different unobserved components or building blocks responsible for the dynamics of the series such as trend, seasonal, cycle, and the effects of explanatory and intervention variables are identified separately before being put together in a state space model.}

% {\color{red} State space methods originated in the eld of control engineering, starting with the groundbreaking
% paper of Kalman (1960). They were initially (and still are) deployed for the purpose
% of accurately tracking the position and velocity of moving objects such as ships, airplanes,
% missiles, and rockets [...] these ideas could well be applied to time series analysis
% generally as well.}

% {\color{red} In a state space
% analysis the time series observations are assumed to depend linearly on a state vector that
% is unobserved and is generated by a stochastically time-varying process (a dynamic system).
% The observations are further assumed to be subject to measurement error that is independent
% of the state vector. The state vector can be estimated or identified once a sufficient set of
% observations becomes available.}

% \begin{definition}
% A linear Gaussian state space model for the multivariate time series $\mathbf{y} = \mathbf{y}_1, \ldots, \mathbf{y_n}$, where each observiation is a $p$ dimensional vector $\mathbf{y}_i = \{ y_{i1}, \ldots, y_{ip} \}$, is given by
% \begin{equation}
% \mathbf{y}_t = \mathbf{Z}_t \mathbf{\alpha}_t + \mathbf{d}_t + \mathbf{\epsilon}_t \quad \mathbf{\epsilon}_t \sim N \left( 0, \mathbf{H}_t \right)
% \end{equation}
% called {\color{red} space?} \emph{orbservation or measurement equation}, and
% \begin{equation}
% \mathbf{\alpha}_t = \mathbf{T}_t \mathbf{\alpha}_{t-1} + \mathbf{c}_t + \mathbf{R}_t \mathbf{\eta}_t \quad \mathbf{\eta}_t \sim N \left( 0, \mathbf{Q}_t \right)
% \end{equation}
% called \emph{state or transition equation}, where the individual summands correspond to:
% \begin{align*}
%  & \mathbf{y}_t \quad \text{observed or measured values,} \\
%  & \mathbf{Z}_t \quad \text{design matrix,} \\
%  & \mathbf{\alpha}_t \quad \text{unobserved state,} \\
%  & \mathbf{d}_t \quad \text{observation intercept,} \\
%  & \mathbf{\epsilon}_t \quad \text{observational disturbance,} \\
%  & \mathbf{H}_t \quad \text{observational disturbance covariance matrix,} \\
%  & \mathbf{T}_t \quad \text{transition matrix,} \\
%  & \mathbf{c}_t \quad \text{state intercept,} \\
%  & \mathbf{R}_t \quad \text{selection matrix,}  \\
%  & \mathbf{\eta}_t \quad \text{state disturbance, and} \\
%  & \mathbf{Q}_t \quad \text{state disturbance covariance matrix}
% \end{align*}
% \end{definition}

% {\color{red} The $p \times m$ matrix $Z_t$ links the observation vector $y_t$ with the unobservable state vector $\alpha_t$ and may consist of regression variables.  The $m \times m$ transition matrix $T_t$ determines the dynamic evolution fo the state vector [...] the observation and state disturbances $\epsilon_t$ and $eta_t$ are assumed to be serially independent and independent of each other at all time points [...] matrix $R_t$ is an $m \times r$ selection matrix with $r < m$.}

% {\color{red} The initial state vector $\alpha_1$ is assumed to be generated as $\alpha_1 \sim NID \left( a_1, P_1 \right)$, independen of the observation and estate disturbances $\epsilon_t$ and $\eta_t$. Mean $a_1$ and variance $P_1$ can be treated as given konw.}

% {\color{red} Talk about initialization?}

% For example, if the time series $\mathbf{y}$ is unidimensional and the state space model is time invariant (only $\mathbf{y}_t$ and $\mathbf{\alpha}_t$ depends on $t$, being the rest of the summands constant), a model with $m$ unobserved stares will be given by
% \[
% y_t = [ z_1 \ldots z_m ]
% \begin{bmatrix}
%   z_{1} \\
%   \vdots \\
%   z_{n}
% \end{bmatrix}
% + d_{t} + \epsilon_{t}
% \]

% Some of the most common time series models are particular cases of the state-space model (see Example XX).

% \begin{example}
% \end{example}

% {\color{red} TODO: Explain the Kalman filter}

% The \emph{Kalman filter} is a recursive formula that provides an optimal estimate for the unknown state in a state space model. At each time step $t$, the Kalman filter computes the predicted state conditional to the observations up to time $t-1$.

% {\color{red} Kalman filter can be use for filtering, prediction and smoothing. Here we are only interested in prediction [...] forward pass [...] recursive formulas}

% \begin{definition}
% \begin{equation}
% \begin{split}
%   & \mathbf{a}_{t+1} = \mathbf{T}_t \mathbf{a}_t + \mathbf{K}_t \mathbf{v}_t  \\
%   & \mathbf{K}_t = \mathbf{T}_t \mathbf{P}_t \mathbf{Z}_t^T \mathbf{F}_t^{1} \\
%   & \mathbf{v}_t = \mathbf{y}_t  - \mathbf{Z}_t \mathbf{a}_t
% \end{split}
% \end{equation}
% called \emph{prediction equations}, and
% \begin{equation}
% \begin{split}
%   & \mathbf{F}_t = \mathbf{Z}_t \mathbf{P}_t \mathbf{Z}_t^T + \mathbf{H}_t \\
%   & \mathbf{L}_t = \mathbf{T}_t - \mathbf{K}_t \mathbf{Z}_t \\
%   & \mathbf{P}_{t+1} = \mathbf{T}_t \mathbf{P}_t \mathbf{L}_t^T + \mathbf{R}_t \mathbf{Q}_t \mathbf{R}_t^T
% \end{split}
% \end{equation}
% called \emph{updating equations}.
% \end{definition}

% \begin{example}
% {\color{red} TODO: Provide an example where we can see how the Kalman filter integrates the predicted probability distribution and observed probability distribution to make a prediction}
% \end{example}

% {\color{red} TODO: Examplain the space search algorithm}

% \begin{example}
% {\color{red} TODO: Provide an example of auto-time series}
% \end{example}

%
% Section: Anomaly Detection
%

\section{Anomaly Detection}

As we have seen in Section XXX, the main problem in anomaly detection is that we do not have a precise mathematical definition of what an anomaly is. In this book we propose to equate \emph{abnormal} samples with \emph{incompressible} samples, and to study the consequences. The essence is that learning is about finding regularities in a dataset, and data compression is about exploiting regularities. We have also seen that the best model minimizes the sum of the length of the model plus the length of the data given the model. This optimal model partitions the dataset into two disjoint subsets: a compressible part and an incompressible part. It is the latter that interests us here: being incompressible means the samples cannot be explained by the model, i.e., they are model-based anomalies under the best available model.

\begin{definition}
Let $\emph{X}$ be a dataset composed by $p$ features and $n$ samples, $\mathbf{y}$ the target variable, and $\mathit{M}$ a model such that the nescience $N \left( \mathbf{X}, \mathit{M} \right)$ is minimal. Let $\hat{y} = \mathit{M} \left( \mathbf{X} \right)$ be the predicitons made by the model $\mathit{M}$ over the vectors of $\mathbf{X}$. We define the \emph{anomaly subset} of $\mathbf{X}$, denoted by $\mathcal{A}_\mathit{M}^\mathbf{X}$, to the set of $X$ such that $y \neq \hat{y}$.
\end{definition}

The class \texttt{nescience.anomalies} allows us to identify the anomaly subset, i.e., the samples that do not match the regularity patterns found in the rest of the dataset. In Example \ref{ex-abnormal_boston_houses} we apply this class to identify houses with abnormally low prices and to explain why they are cheaper.

\begin{example}
\label{ex-abnormal_boston_houses}
We use the Boston House Prices dataset from \texttt{scikit-learn}. The dataset contains 13 predictive features (both numeric and categorical) describing different characteristics of the houses (number of rooms, age, etc.) and the target is the median value of owner-occupied homes. The dataset has 506 samples. We aim to identify houses whose prices are \emph{abnormally low}, i.e., houses that, given their characteristics, should have a higher price.

\begin{sourcecode}
{\scriptsize \begin{verbatim}
from sklearn.datasets import load_boston

data = load_boston()
X    = data.data
y    = data.target
\end{verbatim}}
\end{sourcecode}

We first train a "knowledge model," i.e., the best model that explains the target variable given the predictors, without overfitting. By default, the \texttt{anomalies} class uses the AutoML capabilities provided by the \texttt{nescience} library.

\begin{sourcecode}
{\scriptsize \begin{verbatim}
from nescience.anomalies import Anomalies

model = Anomalies(X_type="mixed", y_type="numeric")
model.fit(X, y)
\end{verbatim}}
\end{sourcecode}

Finally, we select those samples for which the actual price is smaller than the price predicted by the model.

\begin{sourcecode}
{\scriptsize \begin{verbatim}
anomalies = model.get_anomalies("smaller")
X.shape, anomalies.shape
((506, 13), (25,))
\end{verbatim}}
\end{sourcecode}

Thus, there are 25 houses with abnormally low price.

\end{example}

The \texttt{anomalies} class also allows us to group the identified anomalies according to shared characteristics.

Let us examine which attributes best describe those abnormal houses.

\begin{sourcecode}
{\scriptsize \begin{verbatim}
model.get_classes(n_dims=1, an_type="smaller", filter_balancedness=True,
                  filter_redundancy=False, filter_repeated_attrs=False)

Attribute1	Attribute2	Inertia	        N Class 0	N Class 1	Ratio
2	        None	        154.953975	9	        16	        0.36
4	        None   	        0.090593	6	        19	        0.24
5	        None	        5.002437	17	        8	        0.68
7	        None	        20.352117	6	        19	        0.24
8	        None	        77.611111	18	        7	        0.72
9	        None	        41737.11111	7	        18	        0.28
10	        None	        26.577436	13	        12	        0.52
12	        None	        430.294828	18	        7	        0.72
\end{verbatim}}
\end{sourcecode}

According to the inertia, the best attribute to classify the anomalies is attribute 4 (nitric oxides concentration). This attribute divides the abnormal houses into two clusters of sizes 6 and 19. Let us examine how price varies with this dimension.

\begin{sourcecode}
{\scriptsize \begin{verbatim}
from sklearn.linear_model import LinearRegression
import matplotlib.pyplot as plt
        
lr = LinearRegression()
lr.fit(X[:,4].reshape(-1, 1), y)
        
plt.scatter(X[:,4].reshape(-1, 1), y)
plt.plot(X[:,4], lr.intercept_ + lr.coef_ * X[:,4], color="red")
plt.xlabel(data.feature_names[4])
plt.ylabel("Price")
\end{verbatim}}
\end{sourcecode}

\begin{figure}[h]
\centering
\includegraphics[width=0.6\textwidth]{nox_price.png}
\caption{Price as a function of NOX.}
\label{figure:nox_vs_price}
\end{figure}

The regression line suggests that house prices are lower in areas with higher nitric oxides concentration. Let us see how the anomalies split along this dimension.

\begin{sourcecode}
{\scriptsize \begin{verbatim}
class0, class1 = model.get_class_points(attribute1=4, attribute2=None,
                 an_type="smaller")

plt.hist(class0)
plt.hist(class1)
plt.ylabel("Count")
plt.xlabel(data.feature_names[4])
plt.show()
\end{verbatim}}
\end{sourcecode}

\begin{figure}[h]
\centering
\includegraphics[width=0.6\textwidth]{nox_histogram.png}
\caption{Histogram of anomalies NOX.}
\label{figure:nox_histogram}
\end{figure}

The analysis suggests that six of the houses have abnormally low prices because they are located in areas with high pollution levels. We can repeat the same analysis with the other attributes; however, their higher inertia indicates that class separation will be less evident. Note that there may be multiple reasons why a house's price is abnormally low. 

%
% Section: Decision Trees
%

\section{Decision Trees}
\label{sec:decision_trees}

In the last sections we have seen how to use the concepts of miscoding, inaccuracy, surfeit and nescience to evaluate the quality of datasets and models, and to automatically select a family of models and search over its hyperparameters to find the best possible description of a topic. In particular, we have studied in detail the family of binary decision trees. The procedure used in the \texttt{fastautoml} library with trees was a mix between a classical approach (a CART algorithm combined with a cost-complexity pruning), and an evaluation of candidate trees using the minimum nescience principle. In this section we are going to see a new algorithm to derive optimal trees, both for classification and regression problems, that is entirely based on the theory of nescience. The new algorithm, by design, avoids the overfitting of the training dataset without losing accuracy, it does not require the optimization of hyperparameters, thus significantly reducing the training time, and it produces much smaller and more shallow trees than traditional algorithms, facilitating the interpretability of the results.

\subsection{Algorithm Description}
\label{sub:tree_algorithm_description}

The following pseudocode shows the proposed algorithm to build a decision tree given a training dataset $(\mathbf{X}, \mathbf{y})$. The procedure is based on a breadth first traversal of trees combined with a greedy approach. It requires a function called \textsc{bestSplit()} that returns the best split of a given subset of the data into two subsets; and a second function, called \textsc{Nescience()} that provides an estimation of the nescience of the current tree. The algorithm is based on two nested loops: the external \textbf{while} loop keeps a list of the candidate nodes to grow, whereas the internal \textbf{for} loop finds the best node to grow the tree. The latter operation requires to check all possible growing options and select the one that minimizes the nescience. The exit point of the algorithm is when there are no more branches to grow. We keep track of the best nescience achieved during the building process and return the associated tree.

\begin{sourcecode}
\label{algorithm:decision_tree}
{\scriptsize \begin{verbatim}

def BUILD_TREE(data)

    nodesList <- list()
    tree <- BESTSPLIT(data)
    bestNescience <- NESCIENCE(tree)
    nodesList.append(tree)

    while not nodesList.empty()
    
        nescience <- bestNescience
        bestNode <- None
        childNode <- None
        
        for i <- 1, nodesList.length()

            node <- nodesList[i]
            
            node.child <- BESTSPLIT(node.ldata)
            tmp <- NESCIENCE(tree)
            if tmp < nescience
                nescience <- tmp
                bestNode <- i
                
            node.left <- None

        if nescience < bestNescience
            node <- nodeList[bestNode]
            bestNescience <- nescience
            nodesList.append(node.left)
                        
            if not node.left.empty() and not node.right.empty()
                nodesList.remove(bestNode)
    
    return tree
\end{verbatim}}
\end{sourcecode}

The main difference of our algorithm from other decision tree building algorithms is in the way the tree is evaluated. Instead of using only accuracy as most of the algorithms do, in addition, we take into account the complexity of the tree (surfeit) to avoid overfitting, and the quality of the subset of data used during the training process (miscoding).

\subsubsection*{Nescience}
\label{sub:tree_nescience}

The calculation of the nescience implemented in the algorithm is based on a Euclidean distance utility function (see Section \ref{sec:machine_learning_nescience}), because that one was the one that produced the best results in the tests we have performed. For the computation of miscoding and inaccuracy, we use the same techniques that the one used in the \texttt{fastautoml} library, described in Section \ref{sec:machine_learning_miscoding} and Section \ref{sec:machine_learning_inaccuracy} respectively. For the implementation of surfeit, we use the same template to describe trees that was used in the \texttt{DecisionTreeClassifier} of the \texttt{AutoClassifier} class, and that was described in Section \ref{subsec:surfeit_decision_trees}. The only difference is that we also allow equalities in the nodes (\texttt{if [attr] = [thresh]}), something not supported by the \texttt{DecisionTreeClassifier} algorithm of the \texttt{scikit-learn} library.

The generic problem of the instability of inaccuracy due to very short models, also applies to this algorithm (see Section \ref{sec:machine_learning_surfeit}), and the particular problem of the algorithms to build decision trees, in which the best local split might not be that one that minimizes the error (see Section \ref{sec:sec:machine_learning_classification}) is also relevant in this case.

The concept of nescience is used in two different ways in our algorithm. For every iteration of the \texttt{for} loop we have to decide which one of the candidate branches of the tree we should develop. Recall that the order in which we develop the branches is important, since it might happens that one branch does not get develped because that would mean increase the sufeit without a sufficiently large decrease of the inaccuracy. The second place is a the end of the \texttt{while} loop, we we keep trac of the nescience of the different building steps, to decide at the end of the algorithm with wich tree we return.

We treat regression problems as classification problems in which we discretizes the continuous target variable $\mathbf{y}$ into $n$ intervals given the number of samples, and using a uniform discretization (see Section \label{sec:codes_continuous_data}). Once the target variable has been discretized, we train a regular classification tree.

\subsubsection*{Splitting Criteria}
\label{sub:tree_splitting_criteria}

Given a subset $\mathbf{Q} \subseteq \mathbf{X}$ we have to find an split for $\mathbf{Q}$ such that the values of $\textbf{y}$ are grouped together. Recall that a split is a pair $\theta = (j,w)$, were $1 \leq j \leq p$ is a feature index and $w$ is the partition point (see Section \ref{subsec:learning_decision_trees}). A split divides the set $\mathbf{Q}$ into two disjoint subsets $\mathbf{Q}_l$ and $\mathbf{Q}_r = \mathbf{Q} \backslash \mathbf{Q}_l$. In case of a continuous variable we have that $\mathbf{Q}_l = \{\mathbf{x}_i \in \mathbf{Q} : x_{ij} \leq w\}$, and if the feature is categorical we define $\mathbf{Q}_l = \{\mathbf{x}_i \in \mathbf{Q} : x_{ij} = w\}$\footnote{Ideally, for the categorical case, instead of a single feature $w$ we should search over all the elements of the power set of the set of features $\mathcal{P} \{1, 2, \ldots, p \}$. Unfortunately, that would imply to check $2^p$ cases, something that is time-expensive from the computational point of view.}.

In Section \ref{subsec:learning_decision_trees} we saw that a common splitting criteria used in practice is to minimize the weighted entropy $\tilde{H}$ of the subsets $\mathbf{Q}_l$ and $\mathbf{Q}_r$, that is, to find an split that it is minimal $\theta^\star = \argmin_\theta \tilde{H}(\mathbf{Q},\theta)$. More explicitly, if $\mathbf{y}$ is a target vector taking values from a set of $k$ labels $\mathcal{G} = \{g_1, \ldots, g_k \}$ (either because is a categorical target or a continuous target that has been discretized into $k$ intervals), and denoting the subsets of $\mathbf{y}$ as $\mathbf{y}^l = \{y_i : \mathbf{x}_i \in \mathbf{Q}_l \}$ and $\mathbf{y}^r = \{y_i : \mathbf{x}_i \in \mathbf{Q}_r \}$, and $n_l$ and $n_r$ are the number of elements of $\mathbf{y}^l$ and $\mathbf{y}^r$ respectively, we have that
\begin{multline}
\tilde{H}(\mathbf{Q},\theta) = \frac{n_l}{n} \left( - \sum_{i=1}^k \frac{ \sum_{j=1}^{n_l} I(y^l_j = g_i)} {n_l} \log_2{ \frac{ \sum_{j=1}^{n_l} I(y^l_j = g_i)} {n_l} } \right) \\
+ \frac{n_r}{n} \left( - \sum_{i=1}^k \frac{ \sum_{j=1}^{n_r} I(y^r_j = g_i)} {n_r} \log_2{ \frac{ \sum_{j=1}^{n_r} I(y^r_j = g_i)} {n_r} } \right)
\end{multline}

In our nescience based decision tree algorithm, the splitting criteria is to minimize the total length of encoding the subsets $\mathbf{Q}_l$ and $\mathbf{Q}_r$ using optimal codes. We have to find the optimal split $ \theta^\star = \argmin_\theta \hat{K}_C(\mathbf{Q} \mid \theta) = \argmin_\theta \{ \hat{K}_{C_l}(\mathbf{Q}_l) + \hat{K}_{C_r}(\mathbf{Q}_r) \}$ where $C_l$ and $C_r$ are the optimal codes given the relative frequencies of the observed values of $\mathbf{y}^l$ and $\mathbf{y}^r$ respectively. The quantity $\hat{K}_C(\mathbf{Q} \mid \theta)$ is computed as:
\begin{multline}
\hat{K}_C(\mathbf{Q} \mid \theta) = \hat{K}_{C_l}(\mathbf{Q}_l) + \hat{K}_{C_r}(\mathbf{Q}_r) = \\
- \sum_{i=1}^k \log_2{ \frac{ \sum_{j=1}^{n_l} I(y^l_j = g_i)} {n_l} } - \sum_{i=1}^k \log_2{ \frac{ \sum_{j=1}^{n_r} I(y^r_j = g_i)} {n_r} }
\end{multline}
I this particular case (if we use as compression algorithm a code with optimal lengths, and continuous variables have been discretized) it turns out that both expressions are equivalent given the following relation:
\[
\tilde{H}(\mathbf{Q},\theta) = \frac{1}{n} \hat{K}_C(\mathbf{Q} \mid \theta)
\]
We prefer to talk of encoding length instead of weighted entropy because it has an easier interpretation in the context of the theory of nescience.

\begin{remark}
Strictly speaking, if we want to implement a decision trees search algorithm fully compliant with the minimum nescience principle, instead of using a total length encoding as splitting criteria, we should had computed the nescience at each split and select that one that makes it minimal. However, early experiments have shown that at local level it works better to group the values of $y$ than to reduce the nescience. Further research is required to confirm and explain this point.
\end{remark}

\subsubsection*{Practical Implementation}

In the web page that accompanies this book\footnote{http://www.mathematicsunknown.com} we provide an open-source implementation of our algorithm in Python. Our software can be used together with other machine learning tools from the \texttt{scikit-learn} library, since we adhere to the their API guidelines. For example, our algorithm can be used as part of an ensemble of classifiers, like the \texttt{BaggingClassifier} meta-estimator, or the results of the classification could be cross-validated with tools like \texttt{cross\_val\_score}. As an example, to provide a model for the breast cancer dataset, we could do something like the following:

\begin{sourcecode}
{\scriptsize \begin{verbatim}
from NescienceDecisionTree import NescienceDecisionTreeClassifier
from sklearn.datasets import load_breast_cancer

data = load_breast_cancer()

model = NescienceDecisionTreeClassifier()
model.fit(data.data, data.target)
print("Score: ", model.score(data.data, data.target))
\end{verbatim}}
\end{sourcecode}

\subsection{Algorithm Evaluation}
\label{sub:algorithm_evaluation}

In this section we are going to evaluate our new algorithm, and compare its performance against the well-known algorithm CART. CART, \emph{Classification and Regression Trees}, is the de-facto standard algorithm used in the machine learning industry for the derivation of decision trees. For this particular experiment we have used the CART implementation provided by \texttt{scikit-learn}.

Figure \ref{figure:data_error_cart} shows a synthetic dataset consisting of $1000$ random points lying on a two dimensional plane, where all the points with an $X1$ attribute less than $50$ are colored blue, and the rest as red. We have artificially introduced a red point, simulating a measurement error, in the blue area. The black lines correspond to the decisions performed by CART. Since the CART algorithm will not stop until all the points have been properly classified, we have to specify an expected count condition to limit the number of splits. The figure correspond to the tree generated by CART setting the \texttt{min\_samples\_leaf} hyperparameter to $5$.

\begin{figure}[h]
\centering
\includegraphics[width=0.6\textwidth]{data_error_cart.jpg}
\caption{Synthetic dataset with CART algorithm splits.}
\label{figure:data_error_cart}
\end{figure}

The tree obtained by applying our algorithm to the dataset of Figigure \ref{figure:data_error_cart} can be seen in Figure \ref{figure:data_error_nes}. The nescience based algorithm does not try to model the error point, since the gain due to an increment in the accuracy does not compensate the surfeit introduced in the model. Recall that the algorithm stops when the total nescience of the tree, based on the measures of miscoding, inaccuracy and surfeit, does not decrease when adding new nodes to the tree. Our algorithm presents a lower sensitivity to the errors found in datasets, at least if the number of errors is small compared with the number of valid points.

\begin{figure}[h]
\centering
\includegraphics[width=0.4\textwidth]{small_error.jpg}
\caption{Decision tree obtained by the nescience algorithm.}
\label{figure:data_error_nes}
\end{figure}

Our second experiment, again with synthetic dataset, is depicted in Figure \ref{figure:blobs}. There, we create two isotropic Gaussian blobs that partially overlap. We start with a standard deviation of $2.5$ for each cluster, so they are easy to separate, and we increase the standard deviation in increments of $0.01$, until we reach $4.5$, which causes significant overlaps. For each value of the standard deviation, we run the experiment $100$ times and we compute the average accuracy for the two algorithms using different datasets for training (70\% of the data) and testing (30\% of the data). The results of this experiment are shown in Figure \ref{figure:accuracy_blobs}.

\begin{figure}[h]
\centering
\includegraphics[width=0.6\textwidth]{blobs.png}
\caption{Isotropic Gaussian Blobs.}
\label{figure:blobs}
\end{figure}

As we can see, the performance of both algorithms, in terms of accuracy, is similar. However we should note that the hyperparameter \texttt{minimum\_leaf\_size} of the CART algorithm has been optimized to achieve the best accuracy. For this particular experiment, the best value was achieved with a minimum leaf size of 26 points. By definition, given the fact that CART has one degree of freedom more than the nescicence algorithm, it should produce better accuracy; something that it is not observed (both algorithms have a mean accuracy of $0.87$.

\begin{figure}[h]
\centering
\includegraphics[width=0.6\textwidth]{accuracy_blobs.png}
\caption{Accuracy of Isotropic Gaussian Blobs.}
\label{figure:accuracy_blobs}
\end{figure}

For each iteration of the experiment, we have also computed the average number of nodes, including internal and leaf nodes, required by the models to properly classify the clouds in the dataset. The results of this measurement are show in Figure \ref{figure:length_nodes}. Our algorithm requires an average of $4$ nodes compared to $23$ nodes for the CART algorithm. Moreover, our algorithm is more stable than CART, in the sense that it produces models of similar complexity when it gets similar input datasets (a standard deviation of $0.31$ compared to $3.77$ for CART).

\begin{figure}[h]
\centering
\includegraphics[width=0.6\textwidth]{nodes_blobs.png}
\caption{Number of Nodes.}
\label{figure:length_nodes}
\end{figure}

In Figure \ref{figure:blob_max_depth} we show the maximum depth of the tree, defined as the longest path from the root of the tree to any of its leaves. The maximum depth of the tree is a good measure of the average time it will require for the model to provide a classification. The nescience algorithm has an average depth of $1.6$ nodes, whereas the average depth yielded by the CART algorithm is $4.8$ nodes.

\begin{figure}[h]
\centering
\includegraphics[width=0.6\textwidth]{max_depth.png}
\caption{Maximum depth of the model.}
\label{figure:blob_max_depth}
\end{figure}

The last part of the evaluation consists in comparing the performance of our algorithm and CART with a collection real datasets. More specifically, we have selected $12$ well known datasets from the UCI Machine Learning Repository. The selected datasets are: diagnosis of breast cancer (\texttt{cancer}), optical recognition of handwritten digits (\texttt{digits}), predicting protein localization sites in gram-negative bacteria (\texttt{yeast}), classification of NASA space shuttle data (\texttt{shuttle}), classification of blocks in web pages (\texttt{page}), segmentation of outdoor images (\texttt{image}), predicting the age of abalones from physical measurements (\texttt{abalone}), predicting the quality of red and white variants of Portuguese wine (\texttt{wine}), filter spam emails (\texttt{spam}), wall-following robot navigation (\texttt{wall}), classification of land use based on Landsat satellite images (\texttt{landsat}), and distinguishing signals from background noise in the MAGIC gamma telescope images (\texttt{magic}). For each dataset, we have repeated the experiment 100 times, by randomly selecting the training (70\%) and testing (30\%) subsets at each iteration.

\begin{figure}[h]
\centering
\includegraphics[width=0.6\textwidth]{cart_nescience_accuracy.png}
\caption{Maximum depth of the model.}
\label{figure:cart_nescience_accuracy}
\end{figure}

In Figure \ref{figure:cart_nescience_accuracy} we compare the accuracy of the resulting models obtained by applying the CART algorithm and the nescience algorithm to the above datasets. In $4$ of the $12$ datasets, our algorithm provides better accuracy than CART. In the remaining $8$ cases, the accuracy is, on average, less than $1\%$ smaller.

\begin{figure}[h]
\centering
\includegraphics[width=0.6\textwidth]{cart_nescience_nodes.png}
\caption{Maximum depth of the model.}
\label{figure:cart_nescience_nodes}
\end{figure}

In Figure \ref{figure:cart_nescience_nodes} it is shown a comparison of the total number nodes (internal nodes plus leaf nodes) of the resulting models. Only for one of the datasets (\texttt{digits}), our model produces a slightly more complex tree that those generated by CART. In the rest of the cases, the number of nodes in the trees generated by the nescience algorithm have between two and three orders of magnitude fewer nodes (in this figure the $y$ axis is in logarithmic scale).

\begin{figure}[h]
\centering
\includegraphics[width=0.6\textwidth]{cart_nescience_depth.png}
\caption{Maximum depth of the model.}
\label{figure:cart_nescience_depth}
\end{figure}

Finally, in Figure \ref{figure:cart_nescience_depth} we provide a comparison of the depth of the tree of the resulting models. Our algorithm always yields a shallower tree than the CART algorithm.

We would like to mention that the nescience algorithm is hihgly robust with respect to the compressor selected or the nescience function implemented. In Table \ref{table:nescience_functions}, we have apply the nescience algorithm to the datasets described above, and evaluate different alternatives for the definition of the nescience function $N (X , M )$: arithmetic mean
\[
N (X , M ) = (\mu(M , D) + \iota(X , M ) + \sigma(M, D))/3,
\]
geometric mean
\[
N (X , M ) = (\mu(M , D) + \iota(X , M ) + \sigma(M, D)) 1/3,
\]
harmonic mean 
\[
N (X , M ) = 3 / (\mu(M , D) + \iota(X , M ) + \sigma(M, D)) -1 ),
\]
Euclidean distance
\[
N (X , M ) = (\mu(M , D) + \iota(X , M ) + \sigma(M, D)) 1/3,
\]
sum
\[
N (X , M ) = \mu(M , D) + \iota(X , M ) + \sigma(M, D),
\]
and product
\[
\mu(M , D) + \iota(X , M ) + \sigma(M, D).
\]
The table shows limited difference between the different functions. 

\begin{table}[h]
\label{table:nescience_functions}
\centering
\begin{tabular}{l l l l l l l l}
\toprule
 & \small{\textbf{Euclid}} & \small{\textbf{Arithm.}} & \small{\textbf{Geometric}} & \small{\textbf{Product}} & \small{\textbf{Addition}} & \small{\textbf{Harmonic}} \\
\midrule
Acc. & 0.758 & 0.784 & 0.803 & 0.803 & 0.784 & 0.81 \\
Stdev. & 0.051 & 0.041 & 0.033 & 0.033 & 0.041 & 0.038 \\
\bottomrule
\end{tabular}
\caption{Comparison of nescience functions}
\end{table}

Similarly, Table \ref{table:compressor} shows the performance of our algorithm when using the LZMA, zlib, and bz2 compressors. We observe that all of them yield similar performance. The above results suggest that the performance our algorithm is independent of the specific choice made for either implementation aspect.

\begin{table}[h]
\label{table:compressor}
\centering
\begin{tabular}{l l l l}
\toprule
 & \textbf{bz2} & \textbf{lzma} & \textbf{zlib} \\
\midrule
Accuracy & 0.813 & 0.804 & 0.81 \\
Stdev & 0.03 & 0.045 & 0.038 \\
\bottomrule
\end{tabular}
\caption{Comparison of compressors}
\end{table}

We emphasize that the CART algorithm requires to optimize a configuration hyperparameter in order to obtain good results, whereas the algorithm proposed in this book does not require from this optimization.

Shallower trees means faster forecasting times when the models used in production, since the number of \texttt{if-else} conditions to be evaluated is smaller. Moreover, smaller trees makes easier to interpret the results by human analysts. and much shorter training times, something very relevant in case of training ensembles of trees, like random forest or boosted trees (although the use of ensembles of models is highly discouraged by the theory of nescience, given their high surfeit).

%
% Section: Algebraic Model Selection
%

\section{Algebraic Model Selection}

As it was the case for the definition of nescience based on the encyclopedic description of research topics, the nescience of structured datasets can be used to evaluate alternative descriptions of research topics (mathematical models), and to identify how far these descriptions are from an ideal perfect knowledge. This evaluation could be used to identify those topics which require further research. Moreover, the same methodology could be applied to collections of datasets to identify our current knowledge of research areas (collections of topics).

If we combine the concept of nescience of a model, with our concepts of relevance and applicability of research topics, we could apply our methodology for the assisted discovery of interesting questions to collections of datasets; a very useful methodology now that big datasets are becoming widely available.

In order to evaluate the methodology developed, we are going to apply it to a particular research topic: \emph{Multipath Wave Propagation and Fading}. The problem at hand is to understand the effect of a propagation environment on a radio signal, such as the one used by wireless devices. The signals reaching the receiving antenna could follow multiple paths, due to atmospheric reflection and refraction, and reflection from water and objects such as buildings. The effects of these multiple wave paths include constructive and destructive interference (fading), and phase shifting of the original signal, resulting a highly complex received signal (see Figure \ref{fig:Multipath-Signal-Propagation}).

\begin{figure}[h]
\centering\includegraphics[scale=0.5]{fading}
\caption{\label{fig:Multipath-Signal-Propagation}Multipath Signal Propagation}
\end{figure}

In many circumstances, it is too complicated to describe all reflection, diffraction and scattering processes that determine the different paths the signal will follow. Rather, it is preferable to describe the probability (stochastic model) that the received signal attains a certain value. We are interested in to analyze how well these stochastic models (our current knowledge) are able to describe what happen in reality.

The \emph{Rayleigh fading model} assumes that the magnitude of a signal will vary randomly, or fade, according to a Rayleigh distribution (the radial component of the sum of two uncorrelated Gaussian random variables). The Rayleigh probability density function of the power signal is given by:

\[
P_{\sigma}\left(x\right)=\frac{1}{\sigma}\exp\left[-\frac{x}{\sigma}\right]
\]

where $\sigma$ is the mean of the received signals. Rayleigh fading is viewed as a reasonable model for the effect of heavily built-up urban environments, when there is no dominant propagation along a line of sight between the transmitter and receiver.

The Rice or \emph{Rician distribution} describes the power of the received signal when the target consists in many small scatterers of approximately equal strength, plus one dominant scatterer whose individual received signal equals all that of all the small scatterers combined (there is a dominant line of sight). The probability density function of the power of the received signal is given by:

\[
P(x)=\frac{1}{\bar{\sigma}}\left(1+a^{2}\right)\exp\left[-a^{2}-\frac{x}{\bar{\sigma}}\left(1+a^{2}\right)\right]I_{0}\left[2a\sqrt{\left(1+a^{2}\right)\frac{x}{\bar{\sigma}}}\right]
\]


where $\bar{\sigma}$ is the mean of the received signals, and it is equal to $\bar{\sigma}=\left(1+a^{2}\right)\bar{\sigma}_{R}$, being $a^{2}\bar{\sigma}_{R}$ the power of the dominant scatterer, and $I_{0}$ is the modified zeroth order Bessel function of the first kind.

An experiment (see Figure \ref{fig:Experimental-Set-Up}) was set up to collect a real dataset to analyze. The experiment was run on a $135 m^2$ office full of obstacles (interacting objects). The transmitter was an Odroid C1 Linux computer with a Ralink RT5370 USB Wifi adapter. The receiver was a (fixed in space) Motorla Moto G mobile phone. Data was collected using the Kismet \footnote{https://www.kismetwireless.net/index.shtml} platform (an 802.11 layer2 wireless network detector, sniffer, and intrusion detection system), with some ad hoc, home made, software extensions, mostly for data aggregation. A total of 3,177 samples (power level measured in dBm) were collected during one hour experiment.

\begin{figure}[h]
\centering\includegraphics[scale=0.2]{experiment}
\caption{\label{fig:Experimental-Set-Up}Experimental Set Up}
\end{figure}

Next table summarizes the results of applying the three considered models (uniform, Raylegh and Rice) and the optimal encoding using a Huffman code:

\begin{table}[h]
\centering
\begin{tabular}{l l l}
\toprule
\textbf{Model} & \textbf{LDM} & \textbf{Nescience} \\
\midrule
Uniform & 17,351 & 1.30 \\
Rayleigh & 13,229 & 0.75 \\
Rice & 11,118 & 0.47 \\
Huffman & 7,541 & - \\
\bottomrule
\end{tabular}
\caption{Nescience of Models}
\end{table}

The uniform model, that is, assuming zero knowledge about the topic covered by the dataset, has a nescience of 1.30. This value is a kind of upper level for the nescience associated with that particular topic and dataset; any model with a higher nescience should be classified as zero knowledge model. If we introduce the knowledge that in a environment with multiple obstacles the signal propagation can be described as a Gaussian process (Rayleigh distribution), we are able to decrease our nescience to 0.75, that is, there were a 43\% improvement in our understanding of the topic. If we add the knowledge that there is usually a strongly dominant signal seen at the receiver caused by a line of sight between the antenna and the mobile phone (Rician distribution), the nescience decreases to 0.47, and so, we have achieved an additional 23\% gain in our understanding. Given that numbers we can conclude that the Rayleigh model increases our knowledge with respect to the uniform model, and that the Rice model does so with respect to Rayleigh. However, the nescience of this last model is 0.47. That means that there still patterns in the dataset that are not explained by the Rice model, or what it is equivalent according to our methodology, there is still some knowledge to discover and learn.

The methodology has been applied to a dataset gathered in a single experiment under a controlled environment, since the goal of this Chapter was to provide a methodology to quantify the nescience of structured datasets, not to evaluate models for signal propagation and fading. In order to to conclude that, in general, the Rice model is an improvement over Rayleigh, a more realistic experiment is required, with multiple datasets gathered in real environments.

%
% Section: The Analysis of the Incompressible
%

% \section{The Analysis of the Incompressible}

% As we have said in Chapter {chap:Introduction}, one of the reasons to understand how things work is to understand the cause-effect relation in systems. We are interested in this cause/effect relation in two ways. That is, if we want to see an effect in a system, we want to understand wich causes trigger that effect. Also, and perhaps more interesting, if we have observed an (probably undesired effect) in a system we would like to discover what has caused that effect, so we can fix it, and revert the normal situation.

% We could use the theory of nescience to model, and modify, those uncommon effect, by means of training a model and looking at the incompressible part of the data.


% A model $\mathcal{M}$ for a dataset $\mathcal{D} = (X, y)$ is a compressed version of that dataset, since the length of the dataset given the model $l(\mathcal{D} \mid \mathcal{M})$ is smaller that then length of the original dataset $l(\mathcal{D})$. The model $\mathcal{M}$ is composed by the regularities found in the dataset (subject to the algorithm used and the families of models considered). What it is left, $\mathcal{D} \mid \mathcal{M}$ is the incompressible part of the dataset, that is, those samples that have no regularity at all, or present a regularity that requires a description longer than the length of the raw data.

% In this section we are going to show the practical applications of analyzing what it is left, that is, the incrompressible samples of a dataset. An element that is incompressible represents a very unlikely, or uncommon, situation of the entity being studied. A incompressible element does not necessarily means a problem, since if a problem is sufficiently common, it can be compressed. An incompressible element is something than cannot be explained given the normal behaviour of the system. Of course, all of this is assuming that our dataset has no errors.

% Once we have found a model that has the lowest possible nescience for a dataset, we could separate those elements that have not been compresed, denoted by XX, and fit a second model. We could argue that it does not make any sense to model the incompressible part, since, it is incompressible. However, the incompressible part is incompressible with respect to the original entity under study, that is, the global systen. And in this new case, we are studying a different entity, namely, the uncommon parts of an entity. It might happen that we can find regularities in this new entity.



%
% Section: References
%

\section*{References}
\label{sec:ML_references}

% TODO: Fix this reference
% The quantity $1 - I(\mathbf{x};\mathbf{y}) / \max\left\{ H(\mathbf{x}),H(\mathbf{y})\right\}$ has previously been proposed in~\cite{ferri2009experimental} as a candidate definition of \emph{normalized mutual information}. However, to the best of our knowledge, it has not been used in practice.

A insightful description of the differences between explanation models and predictive models, how these models are used in different scientific disciplines, and what are the implications for the process of statistical modeling can be found in \cite{shmueli2010explain}.

The application of the minimum description length principle to the identification of optimal decision trees have been proposed in \cite{quinlan1989inferring}, further refined and clarified in \cite{wallace1993coding}; however the coding method proposed by those authors is different from the one used in this book.

The Minimum Description Length \cite{grunwald2007minimum} and the Minimum Message Length \cite{wallace2005statistical} techniques have been applied to the problem of inferring decision trees in \cite{quinlan1989inferring}, later on clarified and extended in~\cite{wallace1993coding}, in \cite{mehta1995mdl} as a technique for pruning, and in \cite{rastogi1998public}, among others. Although the underlining concepts behind the cost function proposed in this chapter are the same (namely, that learning is equivalent to the capability to compress), our approach is very different from the ones described in these works.


% {\color{red} TODO: Find the original reference of the AirPassengers dataset.}

% {\color{red} TODO: Find the original reference of the Appliance energy consumption dataset. https://archive.ics.uci.edu/ml/datasets/Appliances+energy+prediction}


% \begin{verbatim}

% The normalized compression distance of two vectors $x$ and $y$ computed using a compressor based on optimal codes is equivalent to the normalized mutual information of these two vectors, that is:
% \[
% NCD_C\left(\mathbf{x},\mathbf{y}\right)=1-\frac{I\left(\mathbf{x};\mathbf{y}\right)}{\max\left\{ H\left(\mathbf{x}\right),H\left(\mathbf{y}\right)\right\} }
% \]
% The quantity $1 - \frac{I\left(\mathbf{x};\mathbf{y}\right)}{\max\left\{ H\left(\mathbf{x}\right),H\left(\mathbf{y}\right)\right\} }$ has been already proposed in~\cite{ferri2009experimental} as a candidate definition of the concept of Normalized Mutual Information. However, to the best of our knowledge, it has not been used in practice.

% \end{verbatim}

The practical approximations to deficiency and surplus implemented in the `mnplib` library have close precedents in information theory. For discrete random variables, the information about a target \(Y\) that is not provided by \(X\) is quantified by the conditional entropy \(H(Y\mid X)\), whereas \(H(X\mid Y)\) quantifies the information contained in \(X\) that is not determined by \(Y\) (Cover and Thomas, 2006). When normalized by the corresponding marginal entropies, these quantities are the complements of the two directional uncertainty coefficients, or Theil's \(U\) (Theil, 1967): \(1-I(X;Y)/H(Y)=H(Y\mid X)/H(Y)\) and \(1-I(X;Y)/H(X)=H(X\mid Y)/H(X)\). The two directions also appear jointly in the *variation of information*, \(H(X\mid Y)+H(Y\mid X)\), introduced as an information-theoretic distance between partitions (Meilă, 2007). A related distinction is present in the information bottleneck principle, which seeks representations that retain information relevant to a target while discarding information from the input that is irrelevant to it (Tishby, Pereira, and Bialek, 1999). The estimators used by `mnplib` provide a practical approximation to the theoretical notions of deficiency and surplus: deficiency approximates the proportion of information required to describe the target that is absent from the available data, while surplus approximates the proportion of information contained in the data that is unnecessary for describing the target. Their separation therefore provides a diagnostic interpretation that mutual information alone does not provide, distinguishing between insufficient relevant information and excessive irrelevant information.

% **References**

% Cover, T. M., and Thomas, J. A. (2006). *Elements of Information Theory*, 2nd ed. Wiley.

% Theil, H. (1967). *Economics and Information Theory*. North-Holland Publishing Company.

% Meilă, M. (2007). “Comparing clusterings—An information based distance.” *Journal of Multivariate Analysis*, 98(5), 873–895.

% Tishby, N., Pereira, F. C., and Bialek, W. (1999). “The Information Bottleneck Method.” *Proceedings of the 37th Annual Allerton Conference on Communication, Control, and Computing*, 368–377.
